
%includo il file che contiene la versione dei documenti
\newcommand{\versioneAnalisiDeiRequisiti}{3.2.0}			
\newcommand{\versioneNormeDiProgetto}{3.2.0}			
\newcommand{\versioneGlossario}{3.2.0}			
\newcommand{\versionePianoDiQualifica}{3.2.0}			
\newcommand{\versionePianoDiProgetto}{3.2.0}	
\newcommand{\versioneStudioDiFattibilita}{1.2.0}
\newcommand{\versioneSpecificaTecnica}{3.2.0}


\newcommand{\Versione}{\versioneSpecificaTecnica{}}	  %Versione Finale
\newcommand{\Data}{2014-01-24}				           %Data di creazione
\newcommand{\DataUltimaModifica}{2014-xx-xx}
\newcommand{\TipoDocumento}{Specifica Tecnica}	       %tipo documento

%includo il file header.tex (logo grande in prima pagina piu qualche altra regola)
%questo file contiene impostazioni comuni per tutte i documenti

%definizione packages utilizzati
\documentclass[a4paper]{article}
\usepackage[utf8x]{inputenc}
\usepackage{enumitem}
\usepackage[italian]{babel}
\usepackage{latexsym}
\usepackage{xparse}
\usepackage{float}
\usepackage{subfloat}
\usepackage{subfig}
\usepackage{fancyhdr}
\usepackage{eurofont}
\usepackage{lastpage}
\usepackage{graphicx}
\usepackage{textcomp}
\usepackage{booktabs}
\usepackage{color}
\usepackage{lscape}
\usepackage{hyperref}
%\usepackage[bookmarks=false]{hyperref}
\hypersetup{colorlinks=true, linkcolor=black, anchorcolor=red, urlcolor=blue}
\usepackage{longtable}
\usepackage{tabularx}
\usepackage{abstract}
\usepackage{appendix}
\usepackage{multicol}
\usepackage{bmpsize}
\usepackage[all]{hypcap}
\usepackage{titlesec}
\usepackage{indentfirst}
\usepackage{lipsum,titletoc}

%\setcounter{secnumdepth}{4}

%****************INIZIO GESTIONE SUBSECTION MULTIPLE
\makeatletter
\newcommand\level[1]{%
\ifcase#1\relax\expandafter\chapter\or
\expandafter\section\or
\expandafter\subsection\or
\expandafter\subsubsection\else
\def\next{\@level{#1}}\expandafter\next
\fi}
\newcommand{\@level}[1]{%
\@startsection{level#1}
{#1}
{\z@}%
{-3.25ex\@plus -1ex \@minus -.2ex}%
{1.5ex \@plus .2ex}%
{\normalfont\normalsize\bfseries}}

\newdimen\@leveldim
\newdimen\@dotsdim
{\normalfont\normalsize
\sbox\z@{0}\global\@leveldim=\wd\z@
\sbox\z@{.}\global\@dotsdim=\wd\z@
}

\newcounter{level4}[subsubsection]
\@namedef{thelevel4}{\thesubsubsection.\arabic{level4}}
\@namedef{level4mark}#1{}
\def\l@section{\@dottedtocline{1}{0pt}{\dimexpr\@leveldim*4+\@dotsdim*1+6pt\relax}}
\def\l@subsection{\@dottedtocline{2}{0pt}{\dimexpr\@leveldim*5+\@dotsdim*2+6pt\relax}}
\def\l@subsubsection{\@dottedtocline{3}{0pt}{\dimexpr\@leveldim*6+\@dotsdim*3+6pt\relax}}
\@namedef{l@level4}{\@dottedtocline{4}{0pt}{\dimexpr\@leveldim*7+\@dotsdim*4+6pt\relax}}

\count@=4
\def\@ncp#1{\number\numexpr\count@+#1\relax}
\loop\ifnum\count@<100
\begingroup\edef\x{\endgroup
\noexpand\newcounter{level\@ncp{1}}[level\number\count@]
\noexpand\@namedef{thelevel\@ncp{1}}{%
\noexpand\@nameuse{thelevel\@ncp{0}}.\noexpand\arabic{level\@ncp{1}}}
\noexpand\@namedef{level\@ncp{1}mark}####1{}%
\noexpand\@namedef{l@level\@ncp{1}}%
{\noexpand\@dottedtocline{\@ncp{1}}{0pt}{\the\dimexpr\@leveldim*\@ncp{5}+\@dotsdim*\@ncp{0}\relax}}}%
\x
\advance\count@\@ne
\repeat
\makeatother
\setcounter{secnumdepth}{100}
\setcounter{tocdepth}{100}
%****************FINE GESTIONE SUBSECTION MULTIPLE

%impostazioni relative alla visualizzazione delle section
%nell'indice
\titlecontents{section}
[0pt]%left indent
{\bfseries}
{\contentslabel{2.3em}}
{\hspace*{-2.3em}}
{\hfill\contentspage}
[]%separator


\oddsidemargin=.15in
\evensidemargin=.15in
\textwidth=6in
\topmargin=-.5in
\parindent=0in
\headheight=1in
\DeclareMathSizes{10}{10}{10}{10} %per piano qualifica
\pagestyle{fancy}
\lhead{
\bfseries {\Large \TipoDocumento}\\
\bfseries Versione: \Versione\\
}
\chead{}
\lhead{
\includegraphics[scale=0.455]{../Logo&Header/apertureHead.png}
}
\lfoot{\bfseries \TipoDocumento{} v\Versione}
\cfoot{}
\rfoot{\thepage\ of \mypageref{LastPage}}
\newcommand{\mypageref}[1]{
\hypersetup{linkcolor=black}\pageref{#1}\hypersetup{linkcolor=black}}
%\userpackage{lipsum}
\renewcommand{\footrulewidth}{0.4pt}

%definizioni comandi comuni utilizzati
\newcommand{\numref}[1]{\textsl{\nameref{#1} (\ref{#1})}}
\newcommand{\NomeGruppo}{Aperture Software}
\newcommand{\Progetto}{MaaP: MongoDB as an admin Platform}
\newcommand{\Prop}{CoffeeStrap}

%definizione tecnologie
\newcommand{\Node}{Node.js}
\newcommand{\NodeJS}{Node.js}
\newcommand{\Nodejs}{Node.js}

\newcommand{\mongodb}{MongoDB}

%tanti sub quanti ne vogliamo! :)
\newcommand{\subsubsubsection}{\level{4}}
\newcommand{\subsubsubsubsection}{\level{5}}
\newcommand{\subsubsubsubsubsection}{\level{6}}
\newcommand{\subsubsubsubsubsubsection}{\level{7}}
\newcommand{\subsubsubsubsubsubsubsection}{\level{8}}


%definizione comando per parola glossario
\newcommand{\gloss}[1]{\emph{#1}\ped{\emph{\tiny{G}}}}

\newcommand{\grassetto}{\textbf}

%per inserire immagini
\newcommand{\immagine}[2]{
\begin{center}
\begin{figure}[H]
\includegraphics[width=\textwidth]{{{#1}}}
\caption{#2}
\label{#1}
\end{figure}
\end{center}
}

\newcommand{\Glossario}{
Al fine di evitare ogni ambiguità nella comprensione del linguaggio utilizzato nel presente documento e, in generale, nella documentazione fornita dal gruppo \NomeGruppo{}, ogni termine tecnico, di difficile comprensione o di necessario approfondimento verrà inserito nel documento \emph{Glossario\_{}v\versioneGlossario{}.pdf}.\\
Saranno in esso definiti e descritti tutti i termini in corsivo e allo stesso tempo marcati da una lettera "G" maiuscola in pedice nella documentazione fornita.
}

\newcommand{\Prodotto}{
Lo scopo del prodotto è produrre un \gloss{framework} per generare interfacce web di amministrazione dei dati di business basato sullo stack \gloss{\Nodejs{}} e \gloss{\mongodb{}}.\\
L'obiettivo è quello di semplificare il lavoro allo \gloss{sviluppatore} che dovrà rispondere in modo rapido e standard alle richieste degli esperti di \gloss{business}.
}

%inizio pagina del documento
\begin{document}
\thispagestyle{empty}

\begin{center}\centerline{
%inserisco il logo grande della prima pagina
\includegraphics[scale=0.8]{../Logo&Header/logo.png}}

%metto il link dell'email sotto al logo
%{\href{mailto:ApertureSWE@gmail.com}{\color[rgb]{0.39,0.37,0.38}%ApertureSWE@gmail.com}}\\ [3pc]

\vspace{0.5in}

%titolo del progetto
{\Huge {\Progetto}}\\[.5pc]

\underline{\hspace{6in}}\\[8pc]

{\Huge {\TipoDocumento}}\\[1pc]
%{\emph{Versione \Versione}}\\
\end{center}

%\vspace{.05in}

%\vspace{.05in}

%informazioni documento
\begin{center}
%\section{Informazioni documento}
\begin{tabular}{r|l}
%\textbf{Nome} &\TipoDocumento \\
\textbf{Versione} & \Versione{} \\
\textbf{Data creazione} & \Data{} \\
\textbf{Data ultima modifica} & \DataUltimaModifica{} \\
\textbf{Stato del Documento} & Formale \\		          %CAMBIARE QUI
\textbf{Uso del Documento} & Esterno \\			          %CAMBIARE QUI
\textbf{Redazione} &  Pinato Giacomo, Mattia Sorgato\\	  %CAMBIARE QUI
& Michele Maso, Fabio Miotto\\
& Alessandro Benetti, Andrea Perin\\
\textbf{Verifica} & Alberto Garbui, Alessandro Benetti\\  %ED ANCHE QUI!
\textbf{Approvazione} & Michele Maso\\				      %CAMBIARE QUI
\textbf{Distribuzione} & \parbox[t]{4cm}{\NomeGruppo{}\\Prof. Vardanega Tullio\\Prof. Cardin Riccardo\\ \Prop{} }\\
\end{tabular}
\end{center}

\vspace{0.05in}

%inizio sommario del documento
\begin{abstract}
\begin{center}
Questo documento si propone di presentare la specifica tecnica e architetturale per la realizzazione del prodotto \Progetto{}.
\end{center}
\end{abstract}

%\vspace{.4in}

%seconda pagina, diario delle modifiche
\newpage
Diario delle modifiche
\begin{center}
\begin{longtable}{|c|c|c|p{0.5\linewidth}|}
\toprule
\textbf{Versione} & \textbf{Data} & \textbf{Autore} & \textbf{Modifiche effettuate}\\

%aggiungere qui una midrule per aggiungere una nuova riga alla tabella
\midrule
4.2.0 & 2014-xx-xx &   (RE) & Approvazione documento.\\
\midrule
4.1.1 & 2014-xx-xx &   (VR) & Verifica documento.\\
\midrule
4.1.0 & 2014-xx-xx &   (VR) & Verifica documento.\\

\midrule
4.0.2 & 2014-04-xx &   (PT) &  .\\
\midrule
4.0.1 & 2014-04-09 & Alberto Garbui (PT) & Effettuate correzioni segnalate dal Committente.\\


\midrule
3.2.0 & 2014-03-24 & Michele Maso (RE) & Approvazione documento.\\
\midrule
3.1.1 & 2014-03-22 & Alessandro Benetti (VR) & Verifica documento.\\
\midrule
3.1.0 & 2014-03-20 & Alberto Garbui (VR) & Verifica documento.\\
\midrule
3.0.10 & 2014-03-17 & Alessandro Benetti (PT) & Stesura tracciamento.\\
\midrule
3.0.9 & 2014-03-14 & Giacomo Pinato (PT) & Diagrammi di sequenza.\\
\midrule
3.0.8 & 2014-03-13 & Giacomo Pinato (PT) & Diagrammi di attività.\\
\midrule
3.0.7 & 2014-03-04 & Andrea Perin (PT) & Descrizione design pattern.\\
\midrule
3.0.6 & 2014-02-18 & Michele Maso (PT) & Componenti e Classi.\\
\midrule
3.0.5 & 2014-02-12 & Fabio Miotto (PT) & Componenti e Classi.\\
\midrule
3.0.4 & 2014-02-05 & Mattia Sorgato (PT) & Descrizione Architettura.\\
\midrule
3.0.3 & 2014-01-29 & Giacomo Pinato (PT) & Tecnologie Utilizzate.\\
\midrule
3.0.2 & 2014-01-27 & Fabio Miotto (AM) & Tecnologie Utilizzate.\\
\midrule
3.0.1 & 2014-01-24 & Giacomo Pinato (PT) & Prima stesura del documento.\\

\bottomrule
\caption{Registro delle modifiche}
\label{tab:changelog}

\end{longtable}
\end{center}

%terza pagina Indice (viene aggiornato in automatico con due compilazioni)
\newpage
\tableofcontents

%pagine successive hanno la lista di tabelle e lista delle figure
%(vengono aggiornate in automatico)
\newpage
\listoftables	%tabelle
\listoffigures 	%elenco delle immagini

%qui inizia la prima pagina ufficiale
\newpage
\section{Introduzione}
\subsection{Scopo del documento}
Il presente documento ha lo scopo di definire la progettazione ad alto livello del progetto \textbf{\gloss{MaaP}}, a partire dai requisiti individuati durante l'Analisi. Verrà presentata l'\gloss{architettura} generale secondo la quale saranno organizzate le varie componenti \gloss{software}, i \gloss{Design Pattern} e le tecnologie utilizzate per poi descrivere più dettagliatamente le varie componenti e relative dipendenze.

\subsection{Scopo del prodotto}
\Prodotto{}

\subsection{Glossario}
\Glossario{}

\subsection{Riferimenti}

\subsubsection{Normativi}
\begin{itemize}
\item \grassetto{Analisi dei requisiti}: \emph{Analisi\_{}dei\_{}Requisiti\_{}v\versioneAnalisiDeiRequisiti{}.pdf};
\item \grassetto{Norme di Progetto}: \emph{Norme\_{}di\_{}Progetto\_{}v\versioneNormeDiProgetto{}.pdf};
\item \grassetto{Verbale}: \emph{Verbale\_esterno\_20140305\_v1.2.0.pdf}.
\end{itemize}

\subsubsection{Informativi}
\begin{itemize}
\item \grassetto{Learning Node}: O'Reilly Shelley Powers;
\item \grassetto{AngularJS}: O'Reilly Brad Green e Shyam Seshadri;
\item \grassetto{Software Engineering (8th edition)}: Ian Sommerville, Pearson Education | Addison-Wesley;
\item \grassetto{Design Patterns}: E. Gamma, R. Helm, R. Johnson, J. Vlissides, Pearson Education | Addison-Wesley;
\item \grassetto{Dall'idea al codice con UML 2}: L. Baresi, L. Lavazza, M. Pianciamore, Pearson Education.
\end{itemize}

\newpage
\section{Tecnologie utilizzate}
In questa sezione verranno elencate e descritte le tecnologie che si utilizzeranno durante lo sviluppo del progetto. 
In particolare la colonna portante del progetto sarà lo \gloss{stack} \gloss{MEAN}, ovvero \gloss{MongoDB}, \gloss{Express}, 
\gloss{AngularJS} e \gloss{Node.js}.

\subsection{MongoDB}
Il \gloss{database} con il quale la nostra applicazione dovrà interagire è realizzato con MongoDB, come specificato nel capitolato. 
MongoDB è la piattaforma di supporto per i database di MaaP. La scelta di Mongo è voluta anche per la sua facile integrazione con le
tecnologie back end di Node.js, grazie ai relativi moduli. Questa tecnologia offre inoltre i seguenti vantaggi:
\begin{itemize}
\item Facile indicizzazione: Ogni campo in MongoDB può diventare un indice;
\item Bilanciamento di carico: MongoDB \gloss{scala} orizzontalmente molto facilmente grazie all'utilizzo di \gloss{Shard};
\item Integrazione con Javascript: \gloss{Query} o altre funzioni scritte in Javascript possono essere eseguite direttamente dal database.
\end{itemize}

\subsection{Mongoose}
Mongoose è uno strumento di modellazione di dati provenienti da database MongoDB, sviluppato per lavorare in ambienti asincroni.
Questo modulo di Node.JS serve a dare uno schema definito ai dati provenienti dai database MongoDB di MaaP.

\subsection{Javascript}
Si è deciso di utilizzare Javascript in quanto è il linguaggio su cui si basano tutte le altre tecnologie che andremo ad utilizzare e 
offre quindi una facile integrazione, oltre ad essere un ottimo linguaggio per applicazioni \gloss{web} e \gloss{client} side.

\subsection{Node.js}
Si è deciso di utilizzare il framework Node.js in quanto richiesto nel capitolato e adatto al progetto. 
Le sue caratteristiche più vantaggiose sono:
\begin{itemize}
\item Modello \gloss{Event-driven}: ovvero "programmazione ad eventi", che si basa su un concetto semplice: il flusso del \gloss{programma} non segue un corso specifico ma è guidato dalle azioni dell'utilizzatore;
\item Modello asincrono: grazie a questa caratteristica è possibile ridurre al minimo i tempi di morti in quanto, nell'attesa del completamento di una operazione, si procede con altri flussi logici.
\item Grande scalabilità: Grazie al modo in cui è implementato, Node.js riesce ad essere largamente scalabile con minimo sforzo.
\end{itemize}

\subsection{Express}
Express è un framework minimale e flessibile per la creazione di applicazioni web basate su Node.js. Esso contiene una varietà di funzioni sia per la creazione di applicazioni a singola pagina che multi-pagina oltre che applicazioni ibride.

\subsection{JSON}
Rappresenta il tipo di messaggi con cui client e \gloss{server} si scambiano informazioni. I vantaggi offerti sono:
\begin{itemize}
\item Semplicità: i messaggi \gloss{JSON} sono più corti rispetto ad altri formati di interscambio, e vengono eseguiti più velocemente dal \gloss{parser}. JSON inoltre risulta più semplice e immediato rispetto ad esempio a XML.
\end{itemize}
Svantaggi:
\begin{itemize}
\item Restrittività: JSON è meno restrittivo rispetto ad XML, e questo può permettere di inserire errori nello scambio di messaggi.
\end{itemize}

\subsection{AngularJS}
AngularJS è un framework open-source scritto in Javascript e mantenuto da Google, adatto a sviluppare applicazioni web. Il suo obiettivo principale è fornire alle applicazioni web le funzionalità di MVC, in modo da rendere sia lo sviluppo che il test più semplici.
Le sue caratteristiche più significative sono:
\begin{itemize}
\item \gloss{Two Way Data-Binding}: Una delle caratteristiche principali di AngularJS. Le modifiche apportate al \gloss{model} si rifletto direttamente sugli elementi del \gloss{DOM}, e le modifiche al DOM si ripercuotono automaticamente sul model. Questo alleggerisce molto il \gloss{codice} necessario a controllare e gestire il DOM, automatizzando il \gloss{processo}.
\item Templates: I \gloss{template} \gloss{HTML} sono compilati dal \gloss{browser} nel DOM, il quale costituisce poi l'\gloss{input} per il compilatore AngularJS. Quest'ultimo poi crea il data binding tra il DOM e lo \gloss{scope} dei dati. Uno dei più grandi vantaggi di questa tecnica è che separa presentazione da implementazione, in quanto i template html possono modificati senza alterare il modo in cui sono inseriti i dati.
\item \gloss{Dependency Injection}: AngularJS dispone di uno strumento di dependency injection, che aiuta gli sviluppatori facilitando la creazione, la comprensione e il \gloss{testing} dell'applicazione.
\item \gloss{Directives}: Le directives possono essere usate per definire tag HTML personalizzati che fungono da \gloss{widget}. Possono inoltre essere usate per "decorare" elementi con comportamenti personalizzati o per manipolare attributi del DOM.
\end{itemize}

\subsection{Passport}
Passport è un \gloss{middleware} per Node.js estremamente flessibile e modulare, applicabile facilmente a qualsiasi applicazione web basata su \gloss{Express}. Esso fornisce tutte le funzionalità per implementare un sistema di autenticazione facile e sicuro mediante username e password, oppure mediante login preesistenti derivanti da sistemi di autenticazione per piattaforme social tipo Facebook, Twitter e molti altri.

\subsection{HTML5}
\gloss{HTML5} è un \gloss{linguaggio di markup} per la strutturazione delle \gloss{pagine web}.\\
Nel progetto MaaP è stato scelto di utilizzare HTML5 perché introduce novità finalizzate soprattutto a migliorare il \gloss{disaccoppiamento} tra struttura, definita dal markup, caratteristiche di resa (tipo di carattere, colori, eccetera), definite dalle direttive di stile, e contenuti di una pagina web, definiti dal testo vero e proprio.\\
Inoltre HTML5 prevede il supporto per la memorizzazione locale di grosse quantità di dati scaricati dal web browser, ideale per consentire l'utilizzo di applicazioni web e quindi per il \gloss{framework} MaaP.

\newpage
\section{Descrizione architettura}
\subsection{Metodo e formalismo di specifica}
Si è deciso di procedere utilizzando un approccio \gloss{Top-down} per l'esposizione dell'architettura dell'applicazione, ovvero descrivendo inizialmente le componenti in generale per poi arrivare a trattarle al particolare.
Si descriveranno i \gloss{package} e i componenti per poi dettagliare le singole classi, specificando per ciascuna di esse il tipo, l'obiettivo e la funzionalità. Poi si passerà ad illustrare degli esempi d'uso di Design Pattern (descritti approfonditamente nell'\gloss{Appendice} A) e le tecnologie utilizzate.

\subsection{Architettura generale}
L'architettura del framework segue un modello di architettura in stile Client-Server che prevede la suddivisione dell'applicazione in due parti: la parte client composta dall'\gloss{interfaccia utente} (Client) e la parte server composta dalla \gloss{business} logic (Controller) e alla gestione dei dati persistenti (ModelServer).\\
La parte Client segue il design pattern \gloss{MVVM} utilizzato da AngularJS che gestisce la parte frontend del progetto, esso accede ai dati attraverso chiamate AJAX utilizzando le API fornite dal server.\\
La parte Server utilizza Node.js ed il framework Express per gestire tutte le richieste \gloss{REST} del client. Grazie all'utilizzo di Mongoose, l'applicativo lato server può interagire con il database MongoDB, recuperare i dati necessari e rispondere al client tramite oggetti JSON opportuni.\\ L'utilizzo del middleware Passport permette di gestire l'autenticazione utente, ottenendo così un'architettura di tipo \gloss{RESTlike}.\\
Un esempio di stack MEAN che rispecchia le funzionalità base del framework MaaP è il seguente:
\immagine{./mean}{Esempio di stack MEAN}

I seguenti diagrammi rappresentano l'architettura ad alto livello del framework, indicando i package, le classi e le relazioni che intercorrono tra queste.

\subsubsection{Package}
\immagine{./Diagrammi/GeneralePackage}{Architettura generale del software - vista package}
Nel precedente diagramma sono presenti le relazioni tra i package Client ed il package Server.\\
Vengono inoltre presentati tutti i sotto-package così da facilitare la comprensione dell'intero sistema.

\subsubsection{Classi}
\immagine{./Diagrammi/GeneraleClassi}{Architettura generale del software}
Nel precedente diagramma è presente l'architettura ad alto livello del software e vengono indicate le classi fondamentali per rappresentare le relazioni dell'architettura Client-server. I diagrammi di sequenza relativi allo scambio di segnali, lo scopo ed il contesto di utilizzo sono presenti nella sezione \ref{sequenza}.

\newpage
\section{Componenti e Classi}

\subsection{MaaP}
\subsubsection{Informazioni sul package}
\immagine{./Diagrammi/ComponentiMaaP}{Componenti MaaP}
\subsubsubsection{Descrizione}
Namespace globale per il progetto. Le relazioni tra i package Server e Client identificano il modello di architettura Client-server.
\subsubsubsection{Sotto-componenti}
\begin{itemize}
\item MaaP::Server
\item MaaP::Client
\end{itemize}

\subsection{MaaP::Server}
\subsubsection{Informazioni sul package}
\immagine{./Diagrammi/ComponentiServer}{Componenti Server}
\subsubsubsection{Descrizione}
Package per il \gloss{componente} Server del modello di architettura Client-server.
\subsubsubsection{Sotto-componenti}
\begin{itemize}
\item MaaP::Server::ModelServer
\item MaaP::Server::Controller
\end{itemize}

\subsection{MaaP::Server::ModelServer}
\subsubsection{Informazioni sul package}
\immagine{./Diagrammi/ComponentiModelServer}{Componente MaaP::Server::ModelServer}
\subsubsubsection{Descrizione}
Package ModelServer per il componente Server del modello di architettura Client-server che gestisce i dati persistenti del sistema.

\subsubsubsection{Sottocomponenti}
\begin{itemize}
\item MaaP::Server::ModelServer::DataManager;
\item MaaP::Server::ModelServer::Database;
\item MaaP::Server::ModelServer::DSL.
\end{itemize}

\subsubsection{MaaP::Server::ModelServer::DataManager}
\subsubsubsection{Informazioni sul package}
\immagine{./Diagrammi/ComponentiDataManager}{Componente MaaP::Server::ModelServer::DataManager}
\subsubsubsection{Descrizione}
Componente parte del ModelServer per la gestione dei dati.
\subsubsubsection{Sotto-componenti}
\begin{itemize}
\item MaaP::Server::ModelServer::DataManager::DatabaseAnalysisManager;
\item MaaP::Server::ModelServer::DataManager::DatabaseUserManager;
\item MaaP::Server::ModelServer::DataManager::IndexManager.
\end{itemize}
\subsubsubsection{Classi}

\subsubsubsubsection{JSonComposer}
\grassetto{Nome}\\
MaaP::Server::ModelServer::DataManager::JSonComposer\\
\grassetto{Descrizione}\\
Classe che costruisce un file JSON a partire dalla struttura di una \gloss{Collection}, o di un \gloss{Document}, e dai suoi dati.\\
\grassetto{Utilizzo}\\
Viene utilizzata dai DatabaseManager per costruire il file JSON da inviare al Controller.

\subsubsubsubsection{IDatabaseManager}
\grassetto{Nome}\\
MaaP::Server::ModelServer::DataManager::IDatabaseManager\\
\grassetto{Descrizione}\\
Interfaccia che rappresenta il gestore dei database. Contiene tutte le operazioni che si possono effettuare sul database e l'elaborazione dei dati recuperati da essi.\\
\grassetto{Utilizzo}\\
Viene utilizzata per la gestione delle richieste inoltrate dal Controller.\\
\grassetto{Classi che ereditano}
\begin{itemize}
\item MaaP::Server::ModelServer::DataManager::DatabaseAnalysisManager::DatabaseAnalysisManager;
\item MaaP::Server::ModelServer::DataManager::DatabaseUserManager::DatabaseUserManager.
\end{itemize}

\subsubsubsubsection{IDataRetriever}
\grassetto{Nome}\\
MaaP::Server::ModelServer::DataManager::IDataRetriever\\
\grassetto{Descrizione}\\
Interfaccia attraverso cui i DatabaseManager dialogano con i database. Contiene le operazioni di lettura e scrittura nei database.\\
\grassetto{Utilizzo}\\
Viene utilizzata per recuperare e inserire dati, sui database, su richiesta dei DataManager.\\
\grassetto{Classi che ereditano}
\begin{itemize}
\item MaaP::Server::ModelServer::DataManager::DatabaseAnalysisManager::DataRetrieverAnalysis;
\item MaaP::Server::ModelServer::DataManager::DatabaseUserManager::DataRetrieverUsers.
\end{itemize}

\subsubsubsection{MaaP::Server::ModelServer::DataManager::DatabaseAnalysisManager}
\subsubsubsubsection{Informazioni sul package}
\immagine{./Diagrammi/ComponentiDatabaseAnalysisManager}{Componente MaaP::Server::ModelServer::DataManager::DatabaseAnalysisManager}
\subsubsubsubsection{Descrizione}
Componente parte del DataManager per la gestione dei dati del database di analisi.
\subsubsubsubsection{Classi}

\subsubsubsubsubsection{DatabaseAnalysisManager}
\grassetto{Nome}\\
MaaP::Server::ModelServer::DataManager::DatabaseAnalysisManager::DatabaseAnalysisManager\\
\grassetto{Descrizione}\\
Classe che rappresenta il gestore dei database di analisi. Contiene tutte le operazioni che si possono effettuare sul database di analisi e l'elaborazione dei dati recuperati da essi.\\
\grassetto{Utilizzo}\\
Viene utilizzata per la gestione delle richieste, relative al database di analisi, inoltrate dal Controller.\\
\grassetto{Classi da cui eredita}
\begin{itemize}
\item MaaP::Server::ModelServer::DataManager::IDatabaseManager;
\end{itemize}
\grassetto{Relazioni con altre classi}
\begin{itemize}
\item\grassetto{MaaP::Server::ModelServer::DataManager::JSonComposer}\\
Relazione uscente, utilizza un riferimento a un oggetto di tipo JsonComposer per ottenere il JSON da spedire;
\item\grassetto{MaaP::Server::ModelServer::DSL::CollectionData}\\
Relazione uscente, utilizza un riferimento a un oggetto CollectionData che contiene la struttura di un \gloss{file di descrizione};
\item\grassetto{MaaP::Server::ModelServer::DataManager::DataAnalysisManager::DataRetrieverAnalysis}\\
Relazione uscente, utilizza un riferimento a un oggetto DataRetrieverAnalysis per relazionarsi con il database di analisi;
\item\grassetto{MaaP::Server::ModelServer::DataManager::IndexManager::IndexManager}\\
Relazione uscente, utilizza un riferimento a un oggetto IndexManager per la creazione degli indici;
\item\grassetto{MaaP::Server::Controller::Dispatcher}\\
Relazione entrante, interazioni con le funzionalità del gestore del database di analisi.
\end{itemize}

\subsubsubsubsubsection{DatabaseRetrieverAnalysis}
\grassetto{Nome}\\
MaaP::Server::ModelServer::DataManager::DatabaseAnalysisManager::DatabaseRetrieverAnalysis\\
\grassetto{Descrizione}\\
Classe che rappresenta l'oggetto per interagire con i database.\\
\grassetto{Utilizzo}\\
Viene utilizzata per inserire e leggere dati sui database di analisi e framework.\\
\grassetto{Classi da cui eredita}
\begin{itemize}
\item MaaP::Server::ModelServer::DataManager::IDataRetriever;
\end{itemize}
\grassetto{Relazioni con altre classi}
\begin{itemize}
\item\grassetto{MaaP::Server::ModelServer::DataManager::DatabaseAnalysisManager::DatabaseAnalysisManager}\\
Relazione entrante, interazione con il database;
\item\grassetto{MaaP::Server::ModelServer::Database::MongooseDBAnalysis}\\
Relazione uscente, utilizza un riferimento a un oggetto di tipo MongooseDBAnalysis per creare lo schema dei dati del database di analisi e per interagire con essi;
\item\grassetto{MaaP::Server::ModelServer::Database::MongooseDBFramework}\\
Relazione uscente, utilizza un riferimento a un oggetto di tipo MongooseDBFramework per creare lo schema dei dati del database del framework e per interagire con essi;
\end{itemize}


\subsubsubsection{MaaP::Server::ModelServer::DataManager::DatabaseUserManager}
\subsubsubsubsection{Informazioni sul package}
\immagine{./Diagrammi/ComponentiDatabaseUserManager}{Componente MaaP::Server::ModelServer::DataManager::DatabaseUserManager}
\subsubsubsubsection{Descrizione}
Componente parte del DataManager per la gestione dei dati del database del framework che comprende sia dati utente che impostazioni del sistema.
\subsubsubsubsection{Classi}

\subsubsubsubsubsection{DatabaseUserManager}
\grassetto{Nome}\\
MaaP::Server::ModelServer::DataManager::DatabaseUserManager::DatabaseUserManager\\
\grassetto{Descrizione}\\
Classe che rappresenta il gestore del database del framework. Contiene tutte le operazioni che si possono effettuare sul database del framework e l'elaborazione dei dati recuperati da esso.\\
\grassetto{Utilizzo}\\
Viene utilizzata per la gestione delle richieste relative al database del framework inoltrate dal Controller.\\
\grassetto{Classi da cui eredita}
\begin{itemize}
\item MaaP::Server::ModelServer::DataManager::IDatabaseManager;
\end{itemize}
\grassetto{Relazioni con altre classi}
\begin{itemize}
\item\grassetto{MaaP::Server::ModelServer::DataManager::JSonComposer}\\
Relazione uscente, utilizza un riferimento a un oggetto di tipo JsonComposer per ottenere il JSON da spedire;
\item\grassetto{MaaP::Server::ModelServer::DataManager::DataUserManager::DataRetrieverUsers}\\
Relazione uscente, utilizza un riferimento a un oggetto DataRetrieverUsers per relazionarsi con il database del framework;
\item\grassetto{MaaP::Server::Controller::Dispatcher}\\
Relazione entrante, interazioni con le funzionalità del gestore del database di analisi.
\end{itemize}

\subsubsubsubsubsection{DataRetrieverUsers}
\grassetto{Nome}\\
MaaP::Server::ModelServer::DataManager::DatabaseUserManager::DataRetrieverUsers\\
Classe che rappresenta l'oggetto per interagire con il database del framework.\\
\grassetto{Utilizzo}\\
Viene utilizzata per inserire e leggere dati sul database del framework.\\
\grassetto{Classi da cui eredita}
\begin{itemize}
\item MaaP::Server::ModelServer::DataManager::IDataRetriever;
\end{itemize}
\grassetto{Relazioni con altre classi}
\begin{itemize}
\item\grassetto{MaaP::Server::ModelServer::DataManager::DatabaseUserManager::DatabaseUserManager}\\
Relazione entrante, interazione con il database;
\item\grassetto{MaaP::Server::ModelServer::Database::MongooseDBFramework}\\
Relazione uscente, utilizza un riferimento a un oggetto di tipo MongooseDBFramework per creare lo schema dei dati del database del framework e per interagire con essi.
\end{itemize}



\subsubsubsection{MaaP::Server::ModelServer::DataManager::IndexManager}
\subsubsubsubsection{Informazioni sul package}
\immagine{./Diagrammi/ComponentiIndexManager}{Componente MaaP::Server::ModelServer::DataManager::IndexManager}
\subsubsubsubsection{Descrizione}
Componente parte del DataManager per la creazione e gestione degli indici.
\subsubsubsubsection{Classi}
\subsubsubsubsubsection{IndexManager}
\grassetto{Nome}\\
MaaP::Server::ModelServer::DataManager::IndexManager::IndexManager\\
\grassetto{Descrizione}\\
Classe che rappresenta il gestore degli indici. Contiene tutte le operazioni per la creazione degli indici.\\
\grassetto{Utilizzo}\\
Viene utilizzata per la creazione di indici personalizzati su richiesta del DatabaseAnalysisManager.\\
\grassetto{Relazioni con altre classi}
\begin{itemize}
\item\grassetto{MaaP::Server::ModelServer::DataManager::DatabaseAnalysisManager::DatabaseAnalysisManager}\\
Relazione entrante, interazione con il database;
\item\grassetto{MaaP::Server::ModelServer::DataManager::Database::MongooseDBAnalysis}\\
Relazione uscente, utilizza un riferimento ad un oggetto di tipo MongooseDBAnalysis per creare lo schema dei dati del database di analisi e per interagire con essi;
\item\grassetto{MaaP::Server::ModelServer::DataManager::Database::Query}\\
Relazione uscente debole, utilizza un riferimento ad un oggetto Query per il recupero delle query più utilizzate.
\end{itemize}

\subsubsection{MaaP::Server::ModelServer::Database}
\subsubsubsection{Informazioni sul package}
\immagine{./Diagrammi/ComponentiDatabase}{Componente MaaP::ModelServer::Database}
\subsubsubsection{Descrizione}
Componente parte del ModelServer per la gestione dei dati.
\subsubsubsection{Classi}

\subsubsubsubsection{MongooseDBAnalysis}
\grassetto{Nome}\\
MaaP::Server::ModelServer::Database::MongooseDBAnalysis\\
\grassetto{Descrizione}\\
Classe che rappresenta l'interfaccia di connessione con il database di analisi.\\
\grassetto{Utilizzo}\\
Viene utilizzata per interfacciarsi con il database di analisi fornendo uno schema adeguato.\\
\grassetto{Classi da cui eredita}
\begin{itemize}
\item MaaP::Server::ModelServer::Database::Mongoose;
\end{itemize}
\grassetto{Relazioni con altre classi}
\begin{itemize}
\item\grassetto{MaaP::Server::ModelServer::DataManager::DatabaseAnalysisManager::DataRetrieverAnalysis}\\
Relazione entrante, interazione con il database di analisi;
\item\grassetto{MaaP::Server::ModelServer::DataManager::IndexManager::IndexManager}\\
Relazione entrante, interazione con il database di analisi;
\item\grassetto{MaaP::Server::ModelServer::Database::DBAnalysis}\\
Relazione uscente debole, utilizza un riferimento al database di analisi a cui connettersi.
\end{itemize}

\subsubsubsubsection{DBAnalysis}
\grassetto{Nome}\\
MaaP::Server::ModelServer::Database::DBAnalysis\\
\grassetto{Descrizione}\\
Classe che rappresenta il database di analisi.\\
\grassetto{Utilizzo}\\
Viene utilizzata per contenere i dati di analisi.\\
\grassetto{Relazioni con altre classi}
\begin{itemize}
\item\grassetto{MaaP::Server::ModelServer::Database::MongooseDBAnalysis}\\
Relazione entrante debole, interazione con il database di analisi.
\end{itemize}

\subsubsubsubsection{Mongoose}
\grassetto{Nome}\\
MaaP::Server::ModelServer::Database::Mongoose\\
\grassetto{Descrizione}\\
Interfaccia che permette di dialogare con i database utilizzando Mongoose.\\
\grassetto{Utilizzo}\\
Viene utilizzata per interfacciarsi con i vari database.\\
\grassetto{Classi che ereditano}
\begin{itemize}
\item MaaP::Server::ModelServer::Database::MongooseDBAnalysis;
\item MaaP::Server::ModelServer::Database::MongooseDBFramework.
\end{itemize}

\subsubsubsubsection{MongooseDBFramework}
\grassetto{Nome}\\
MaaP::Server::ModelServer::Database::MongooseDBFramework\\
\grassetto{Descrizione}\\
Classe che rappresenta l'interfaccia di connessione con il database del framework.\\
\grassetto{Utilizzo}\\
Viene utilizzata per interfacciarsi con il database del framework fornendo uno schema adeguato.\\
\grassetto{Classi da cui eredita}
\begin{itemize}
\item MaaP::Server::ModelServer::Database::Mongoose;
\end{itemize}
\grassetto{Relazioni con altre classi}
\begin{itemize}
\item\grassetto{MaaP::Server::ModelServer::DataManager::DatabaseAnalysisManager::DataRetrieverAnalysis}\\
Relazione entrante, interazione con il database del framework;
\item\grassetto{MaaP::Server::ModelServer::DataManager::DatabaseUserManager::DataRetrieverUsers}\\
Relazione entrante, interazione con il database del framework;
\item\grassetto{MaaP::Server::ModelServer::Database::DBFramework}\\
Relazione uscente debole, utilizza un riferimento al database del framework a cui connettersi.
\end{itemize}


\subsubsubsubsection{DBFramework}
\grassetto{Nome}\\
MaaP::Server::ModelServer::Database::DBFramework\\
\grassetto{Descrizione}\\
Classe che rappresenta il database del framework.\\
\grassetto{Utilizzo}\\
Viene utilizzata per contenere i dati utente ed impostazioni varie del sistema.\\
\grassetto{Relazioni con altre classi}
\begin{itemize}
\item\grassetto{MaaP::Server::ModelServer::Database::User}\\
Relazione uscente, utilizza un riferimento ad un oggetto \gloss{User} per gestire i dati utente.
\item\grassetto{MaaP::Server::ModelServer::Database::Query}\\
Relazione uscente, utilizza un riferimento ad un oggetto Query per gestire la lista di query fin'ora effettuate dal sistema;
\item\grassetto{MaaP::Server::Controller::FrontController}\\
Relazione entrante, interazione con il database del framework;
\item\grassetto{MaaP::Server::Controller::Passport}\\
Relazione entrante debole, interazione con il database del framework.
\end{itemize}

\subsubsubsubsection{User}
\grassetto{Nome}\\
MaaP::Server::ModelServer::Database::User\\
\grassetto{Descrizione}\\
Classe che rappresenta la parte contenuta nel database del framework relativa ai dati utenti.\\
\grassetto{Utilizzo}\\
Viene utilizzata per contenere i dati utente.\\
\grassetto{Relazioni con altre classi}
\begin{itemize}
\item\grassetto{MaaP::Server::ModelServer::Database::DBFramework}\\
Relazione entrante, interazione con i dati utente.
\end{itemize}

\subsubsubsubsection{Query}
\grassetto{Nome}\\
MaaP::Server::ModelServer::Database::Query\\
\grassetto{Descrizione}\\
Classe che rappresenta la parte contenuta nel database del framework relativa alle query effettuate del sistema.\\
\grassetto{Utilizzo}\\
Viene utilizzata per contenere le query effettuate del sistema.\\
\grassetto{Relazioni con altre classi}
\begin{itemize}
\item\grassetto{MaaP::Server::ModelServer::Database::DBFramework}\\
Relazione entrante, interazione con le query effettuate del sistema.
\item\grassetto{MaaP::Server::ModelServer::DataManager::IndexManager::IndexManager}\\
Relazione entrante debole, interazione con le query effettuate del sistema.
\end{itemize}


\subsubsection{MaaP::Server::ModelServer::DSL}
\subsubsubsection{Informazioni sul package}
\immagine{./Diagrammi/ComponentiDSL}{Componente MaaP::ModelServer::DSL}
\subsubsubsection{Descrizione}
Componente parte del ServerModel per la gestione dei file di descrizione.
\subsubsubsection{Classi}

\subsubsubsubsection{ParserInterface}
\grassetto{Nome}\\
MaaP::Server::ModelServer::DSL::ParserInterface\\
\grassetto{Descrizione}\\
Interfaccia che rappresenta la componente interfaccia del design pattern \gloss{strategy} per il parser di un linguaggio DSL.\\
\grassetto{Utilizzo}\\
Viene utilizzata per la effettuare il parsing di un file di descrizione.\\
\grassetto{Classi che ereditano}
\begin{itemize}
\item MaaP::Server::ModelServer::DSL::DSLParser.
\end{itemize}

\subsubsubsubsection{DSLParser}
\grassetto{Nome}\\
MaaP::Server::ModelServer::DSL::DSLParser\\
\grassetto{Descrizione}\\
Classe che rappresenta l'\gloss{algoritmo} per il parser DSL del design pattern strategy.\\
\grassetto{Utilizzo}\\
Viene utilizzata all'avvio del sistema per eseguire il parsing dei file di descrizione.\\
\grassetto{Classi da cui eredita}
\begin{itemize}
\item MaaP::Server::ModelServer::DSL::ParserInterface;
\end{itemize}
\grassetto{Relazioni con altre classi}
\begin{itemize}
\item\grassetto{MaaP::Server::ModelServer::DSL::DSLDescriptionFile}\\
Relazione uscente debole, utilizza un riferimento ad un oggetto DSLDescriptionFile per leggere il file di descrizione;
\end{itemize}

\subsubsubsubsection{DSLManager}
\grassetto{Nome}\\
MaaP::Server::ModelServer::DSL::DSLManager\\
\grassetto{Descrizione}\\
Classe che rappresenta il gestore dei file di descrizione. Contiene tutte le operazioni per eseguire il parsing dei file di descrizione e per salvare il risultato su appositi file di tipo CollectionData\\
\grassetto{Utilizzo}\\
Viene utilizzata all'avvio del sistema per eseguire il parsing dei file di descrizione e salvare il risultato su file.\\
\grassetto{Relazioni con altre classi}
\begin{itemize}
\item\grassetto{MaaP::Server::ModelServer::DSL::ParserInterface}\\
Relazione uscente, utilizza un riferimento ad un oggetto ParserInterface per eseguire il parsing del file di descrizione;
\item\grassetto{MaaP::Server::ModelServer::DSL::DSLDescriptionFile}\\
Relazione uscente, utilizza un riferimento ad un oggetto DSLDescriptionFile per leggere il file di descrizione;
\item\grassetto{MaaP::Server::ModelServer::DSL::CollectionData}\\
Relazione uscente, utilizza un riferimento ad un oggetto CollectionData per salvare i risultati dell'operazione di parsing.
\end{itemize}

\subsubsubsubsection{CollectionData}
\grassetto{Nome}\\
MaaP::Server::ModelServer::DSL::CollectionData\\
\grassetto{Descrizione}\\
Classe che rappresenta il file contenente il risultato dell'operazione di parsing.\\
\grassetto{Utilizzo}\\
Viene utilizzata all'avvio del sistema per salvare il risultato dell'operazione di parsing del file di descrizione.\\
\grassetto{Relazioni con altre classi}
\begin{itemize}
\item\grassetto{MaaP::Server::ModelServer::DSL::DSLManager}\\
Relazione entrante, interazione con il file;
\item\grassetto{MaaP::Server::ModelServer::DataManager::DatabaseAnalysisManager::DatabaseAnalysisManager}\\
Relazione entrante, interazione con il file.
\end{itemize}


\subsection{MaaP::Server::Controller}
\subsubsection{Informazioni sul package}
\immagine{./Diagrammi/ComponentiController}{Componente MaaP::Server::Controller}
\subsubsubsection{Descrizione}
Package per il componente Controller del modello di architettura Client-server.
\subsubsubsection{Classi}

\subsubsubsubsection{IPassport}
\grassetto{Nome}\\
MaaP::Server::Controller::IPassport\\
\grassetto{Descrizione}\\
Interfaccia che rappresenta il componente target del design pattern object \gloss{adapter}.\\
\grassetto{Utilizzo}\\
Viene utilizzata per gestire l'autenticazione utente.\\
\grassetto{Relazioni con altre classi}
\begin{itemize}
\item\grassetto{MaaP::Server::Controller::FrontController}\\
Relazione entrante, interazione con il gestore dell'autenticazione.
\end{itemize}
\grassetto{Classi che ereditano}
\begin{itemize}
\item MaaP::Server::Controller::PassportAdapter.
\end{itemize}

\subsubsubsubsection{PassportAdapter}
\grassetto{Nome}\\
MaaP::Server::Controller::PassportAdapter\\
\grassetto{Descrizione}\\
Classe che rappresenta il componente adapter del design pattern object adapter.\\
\grassetto{Utilizzo}\\
Viene utilizzata per gestire l'autenticazione utente.\\
\grassetto{Classi da cui eredita}
\begin{itemize}
\item MaaP::Server::Controller::IPassport.
\end{itemize}
\grassetto{Relazioni con altre classi}
\begin{itemize}
\item\grassetto{MaaP::Server::Controller::Passport}\\
Relazione uscente, utilizza un riferimento ad un oggetto di tipo Passport per gestire l'autenticazione utente.
\end{itemize}

\subsubsubsubsection{Passport}
\grassetto{Nome}\\
MaaP::Server::Controller::Passport\\
\grassetto{Descrizione}\\
Classe che rappresenta il componente adaptee del design pattern object adapter.\\
\grassetto{Utilizzo}\\
Viene utilizzata per gestire l'autenticazione utente.\\
\grassetto{Relazioni con altre classi}
\begin{itemize}
\item\grassetto{MaaP::Server::ModelServer::Database::MongooseDBFramework}\\
Relazione uscente debole, utilizza un riferimento ad un oggetto MongooseDBFramework per accedere ai dati utente.
\end{itemize}

\subsubsubsubsection{FrontController}
\grassetto{Nome}\\
MaaP::Server::Controller::FrontController\\
\grassetto{Descrizione}\\
Classe che rappresenta il componente controller del design pattern Front Controller.\\
\grassetto{Utilizzo}\\
Viene utilizzata per gestire le richieste del client ed inoltrarle al dispatcher.\\
\grassetto{Relazioni con altre classi}
\begin{itemize}
\item\grassetto{MaaP::Server::Controller::IPassport}\\
Relazione uscente, contiene un riferimento ad un oggetto IPassport per gestire l'autenticazione utente;
\item\grassetto{MaaP::Server::Controller::Dispatcher}\\
Relazione uscente, contiene un riferimento ad un oggetto Dispatcher per smistare le richieste del client ai vari manager;
\item\grassetto{MaaP::Server::ModelServer::Database::MongooseDBFramework}\\
Relazione uscente, contiene un riferimento ad un oggetto MongooseDBFramework per inserire nuovi dati nel database del framework relativi a nuovi utenti;
\item\grassetto{MaaP::Client::ModelClient::Services::HTTP}\\
Relazione entrante debole, interazione con il \gloss{servizio} HTTP.
\end{itemize}

\subsubsubsubsection{Dispatcher}
\grassetto{Nome}\\
MaaP::Server::Controller::Dispatcher\\
\grassetto{Descrizione}\\
Classe che rappresenta il componente dispatcher del design pattern Front Controller.\\
\grassetto{Utilizzo}\\
Viene utilizzata per smistare le richieste del client ai vari gestori dei dati.\\
\grassetto{Relazioni con altre classi}
\begin{itemize}
\item\grassetto{MaaP::Server::Controller::FrontController}\\
Relazione entrante, interazione con il FrontController;
\item\grassetto{MaaP::Server::ModelServer::DataManager::DatabaseAnalysisManager::DatabaseAnalysisManager}\\
Relazione uscente, contiene un riferimento ad un oggetto DatabaseAnalysisManager per richiedere azioni relative ai dati di analisi;
\item\grassetto{MaaP::Server::ModelServer::DataManager::DatabaseUserManager::DatabaseUserManager}\\
Relazione uscente, contiene un riferimento ad un oggetto DatabaseUserManager per richiedere azioni relative ai dati utenti ed impostazioni di sistema.
\end{itemize}


\subsection{MaaP::Client}
\subsubsection{Informazioni sul package}
\immagine{./Diagrammi/ComponentiClient}{Componente MaaP::Client}
\subsubsubsection{Descrizione}
Package per il componente Client del modello di architettura Client-server.
\subsubsubsection{Sottocomponenti}
\begin{itemize}
\item MaaP::Client::View;
\item MaaP::Client::ControllerModelView;
\item MaaP::Client::ModelClient.
\end{itemize}

\subsubsection{MaaP::Client::View}
\subsubsubsection{Informazioni sul package}
\immagine{./Diagrammi/ComponentiView}{Componente MaaP::Client::View}
\subsubsubsection{Descrizione}
Componente che contiene i template per la visualizzazione delle pagine web.

\subsubsubsubsection{404}
\grassetto{Nome}\\
MaaP::Client::View::404\\
\grassetto{Descrizione}\\
Classe che rappresenta il template per la pagina d'errore.\\
\grassetto{Utilizzo}\\
Viene utilizzata per renderizzare la pagina web d'errore 404.\\

\subsubsubsubsection{collection}
\grassetto{Nome}\\
MaaP::Client::View::collection\\
\grassetto{Descrizione}\\
Classe che rappresenta il template per la parte della pagina destinata alla visualizzazione di una collection.\\
\grassetto{Utilizzo}\\
Viene utilizzata per visualizzare una serie di document di una determinata collection.\\
\grassetto{Relazioni con altre classi}
\begin{itemize}
\item\grassetto{MaaP::Client::ControllerModelView::ControllerClient::CollectionCtrl}\\
Relazione uscente, contiene un riferimento ad un oggetto CollectionCtrl per gestire la creazione della pagina collection-index.
\end{itemize}

\subsubsubsubsection{dashboard}
\grassetto{Nome}\\
MaaP::Client::View::dashboard\\
\grassetto{Descrizione}\\
Classe che rappresenta il template per la parte della pagina destinata alla visualizzazione della dashboard.\\
\grassetto{Utilizzo}\\
Viene utilizzata per visualizzare una serie di collection presenti nel progetto.\\
\grassetto{Relazioni con altre classi}
\begin{itemize}
\item\grassetto{MaaP::Client::ControllerModelView::ControllerClient::DashboardCtrl}\\
Relazione uscente, contiene un riferimento ad un oggetto DashboardCtrl per gestire la creazione della dashboard.
\end{itemize}

\subsubsubsubsection{document}
\grassetto{Nome}\\
MaaP::Client::View::document\\
\grassetto{Descrizione}\\
Classe che rappresenta il template per la parte della pagina destinata alla visualizzazione di un document.\\
\grassetto{Utilizzo}\\
Viene utilizzata per visualizzare la serie di dati che compongono un document.\\
\grassetto{Relazioni con altre classi}
\begin{itemize}
\item\grassetto{MaaP::Client::ControllerModelView::ControllerClient::DocumentCtrl}\\
Relazione uscente, contiene un riferimento ad un oggetto DocumentCtrl per gestire la creazione della pagina document-show.
\end{itemize}

\subsubsubsubsection{documentEdit}
\grassetto{Nome}\\
MaaP::Client::View::documentEdit\\
\grassetto{Descrizione}\\
Classe che rappresenta il template per la parte della pagina destinata alla modifica di un document.\\
\grassetto{Utilizzo}\\
Viene utilizzata per visualizzare la serie di dati che compongono un document per la successiva modifica.\\
\grassetto{Relazioni con altre classi}
\begin{itemize}
\item\grassetto{MaaP::Client::ControllerModelView::ControllerClient::DocumentEditCtrl}\\
Relazione uscente, contiene un riferimento ad un oggetto DocumentEditCtrl per gestire la creazione della pagina di modifica di un document.
\end{itemize}

\subsubsubsubsection{indexCollection}
\grassetto{Nome}\\
MaaP::Client::View::indexCollection\\
\grassetto{Descrizione}\\
Classe che rappresenta il template per la parte della pagina destinata alla visualizzazione degli indici.\\
\grassetto{Utilizzo}\\
Viene utilizzata per visualizzare la serie di indici attualmente presenti nel database di analisi.\\
\grassetto{Relazioni con altre classi}
\begin{itemize}
\item\grassetto{MaaP::Client::ControllerModelView::ControllerClient::IndexCtrl}\\
Relazione uscente, contiene un riferimento ad un oggetto IndextCtrl per gestire la creazione della pagina di visualizzazione degli indici.
\end{itemize}

\subsubsubsubsection{login}
\grassetto{Nome}\\
MaaP::Client::View::login\\
\grassetto{Descrizione}\\
Classe che rappresenta il template per la parte della pagina destinata al login.\\
\grassetto{Utilizzo}\\
Viene utilizzata per visualizzare la pagina di login.\\
\grassetto{Relazioni con altre classi}
\begin{itemize}
\item\grassetto{MaaP::Client::ControllerModelView::ControllerClient::LoginCtrl}\\
Relazione uscente, contiene un riferimento ad un oggetto LoginCtrl per gestire la creazione della pagina di login.
\end{itemize}

\subsubsubsubsection{logout}
\grassetto{Nome}\\
MaaP::Client::View::logout\\
\grassetto{Descrizione}\\
Classe che rappresenta il template per la parte della pagina destinata al logout.\\
\grassetto{Utilizzo}\\
Viene utilizzata per visualizzare la pagina di logout.\\

\subsubsubsubsection{Navbar}
\grassetto{Nome}\\
MaaP::Client::View::Navbar\\
\grassetto{Descrizione}\\
Classe che rappresenta il template per la parte della pagina che visualizza la barra di navigazione.\\
\grassetto{Utilizzo}\\
Viene utilizzata per visualizzare la barra di navigazione.\\
\grassetto{Relazioni con altre classi}
\begin{itemize}
\item\grassetto{MaaP::Client::ControllerModelView::ControllerClient::NavBarCtrl}\\
Relazione uscente, contiene un riferimento ad un oggetto NavBarCtrl per gestire la creazione della barra di navigazione del sito.
\end{itemize}

\subsubsubsubsection{queryCollection}
\grassetto{Nome}\\
MaaP::Client::View::queryCollection\\
\grassetto{Descrizione}\\
Classe che rappresenta il template per la parte della pagina che visualizza la lista delle query più utilizzate.\\
\grassetto{Utilizzo}\\
Viene utilizzata per visualizzare le query più utilizzate.\\
\grassetto{Relazioni con altre classi}
\begin{itemize}
\item\grassetto{MaaP::Client::ControllerModelView::ControllerClient::QueryCtrl}\\
Relazione uscente, contiene un riferimento ad un oggetto QueryCtrl per gestire la creazione della pagina delle query più utilizzate.
\end{itemize}

\subsubsubsubsection{register}
\grassetto{Nome}\\
MaaP::Client::View::register\\
\grassetto{Descrizione}\\
Classe che rappresenta il template per la parte della pagina di registrazione utente.\\
\grassetto{Utilizzo}\\
Viene utilizzata per visualizzare il form di registrazione utente.\\
\grassetto{Relazioni con altre classi}
\begin{itemize}
\item\grassetto{MaaP::Client::ControllerModelView::ControllerClient::RegisterCtrl}\\
Relazione uscente, contiene un riferimento ad un oggetto RegisterCtrl per gestire la registrazione di un nuovo utente.
\end{itemize}

\subsubsubsubsection{userCollection}
\grassetto{Nome}\\
MaaP::Client::View::userCollection\\
\grassetto{Descrizione}\\
Classe che rappresenta il template per la parte della pagina che visualizza la lista degli utenti registrati nel sistema.\\
\grassetto{Utilizzo}\\
Viene utilizzata per visualizzare la pagina contenente la lista degli utenti registrati nel sistema per l'utente amministratore.\\
\grassetto{Relazioni con altre classi}
\begin{itemize}
\item\grassetto{MaaP::Client::ControllerModelView::ControllerClient::UsersCollectionCtrl}\\
Relazione uscente, contiene un riferimento ad un oggetto UsersCollectionCtrl per gestire la creazione della lista di utenti registrati nel sistema.
\end{itemize}

\subsubsubsubsection{userEdit}
\grassetto{Nome}\\
MaaP::Client::View::userEdit\\
\grassetto{Descrizione}\\
Classe che rappresenta il template per la parte della pagina che modifica il profilo utente.\\
\grassetto{Utilizzo}\\
Viene utilizzata per visualizzare il form di modifica del profilo utente.\\
\grassetto{Relazioni con altre classi}
\begin{itemize}
\item\grassetto{MaaP::Client::ControllerModelView::ControllerClient::UserEditCtrl}\\
Relazione uscente, contiene un riferimento ad un oggetto UserEditCtrl per gestire la modifica di un profilo utente da parte di un utente amministratore.
\item\grassetto{MaaP::Client::ControllerModelView::ControllerClient::ProfileEditCtrl}\\
Relazione uscente, contiene un riferimento ad un oggetto ProfileEditCtrl per gestire la modifica del proprio profilo utente.
\end{itemize}

\subsubsubsubsection{userProfile}
\grassetto{Nome}\\
MaaP::Client::View::userProfile\\
\grassetto{Descrizione}\\
Classe che rappresenta il template per la parte della pagina che visualizza il profilo utente.\\
\grassetto{Utilizzo}\\
Viene utilizzata per visualizzare i dati del profilo utente.\\
\grassetto{Relazioni con altre classi}
\begin{itemize}
\item\grassetto{MaaP::Client::ControllerModelView::ControllerClient::UserCtrl}\\
Relazione uscente, contiene un riferimento ad un oggetto UserCtrl per gestire la visualizzazione di un profilo utente da parte di un utente amministratore.
\item\grassetto{MaaP::Client::ControllerModelView::ControllerClient::ProfileCtrl}\\
Relazione uscente, contiene un riferimento ad un oggetto ProfileCtrl per gestire la visualizzazione del proprio profilo utente.
\end{itemize}


\subsubsection{MaaP::Client::ControllerModelView}
\subsubsubsection{Informazioni sul package}
\immagine{./Diagrammi/ComponentiControllerModelView}{Componente MaaP::Client::ControllerModelView}
\subsubsubsection{Descrizione}
Componente ModelView del design pattern MVVM.
\subsubsubsection{Sotto-componenti}
\begin{itemize}
\item MaaP::Client::ControllerModelView::ControllerClient;
\end{itemize}

\subsubsubsection{MaaP::Client::ControllerModelView::ControllerClient}
\subsubsubsubsection{Informazioni sul package}
\immagine{./Diagrammi/ComponentiControllerClient}{Componente MaaP::Client::ControllerModelView::ControllerClient}
\subsubsubsubsection{Descrizione}
Componente parte del ControllerModelView contenente i vari controller.
\subsubsubsubsection{Classi}

\subsubsubsubsubsection{CollectionCtrl}
\grassetto{Nome}\\
MaaP::Client::ControllerModelView::ControllerClient::CollectionCtrl\\
\grassetto{Descrizione}\\
Classe che rappresenta il controller per indirizzare le richieste di visualizzazione di una collection.\\
\grassetto{Utilizzo}\\
Viene utilizzata per la indirizzare le richieste di visualizzazione di una collection.\\
\grassetto{Relazioni con altre classi}
\begin{itemize}
\item\grassetto{MaaP::Client::ModelClient::Services::CollectionDataService}\\
Relazione uscente debole, contiene un riferimento ad un oggetto CollectionDataService per utilizzare il relativo servizio;
\item\grassetto{MaaP::Client::ModelClient::Services::DocumentEditService}\\
Relazione uscente debole, contiene un riferimento ad un oggetto DocumentEditService per utilizzare il relativo servizio;
\item\grassetto{MaaP::Client::View::collection}\\
Relazione entrante, interazione con il template.
\end{itemize}

\subsubsubsubsubsection{DashboardCtrl}
\grassetto{Nome}\\
MaaP::Client::ControllerModelView::ControllerClient::DashboardCtrl\\
\grassetto{Descrizione}\\
Classe che rappresenta il controller per indirizzare le richieste di visualizzazione della dashboard.\\
\grassetto{Utilizzo}\\
Viene utilizzata per la indirizzare le richieste di visualizzazione della dashboard.\\
\grassetto{Relazioni con altre classi}
\begin{itemize}
\item\grassetto{MaaP::Client::ModelClient::Services::CollectionListService}\\
Relazione uscente debole, contiene un riferimento ad un oggetto CollectionListService per utilizzare il relativo servizio;
\item\grassetto{MaaP::Client::View::dashboard}\\
Relazione entrante, interazione con il template.
\end{itemize}

\subsubsubsubsubsection{DocumentCtrl}
\grassetto{Nome}\\
MaaP::Client::ControllerModelView::ControllerClient::DocumentCtrl\\
\grassetto{Descrizione}\\
Classe che rappresenta il controller per indirizzare le richieste di visualizzazione di un document.\\
\grassetto{Utilizzo}\\
Viene utilizzata per la indirizzare le richieste di visualizzazione della pagina document-show.\\
\grassetto{Relazioni con altre classi}
\begin{itemize}
\item\grassetto{MaaP::Client::ModelClient::Services::DocumentDataService}\\
Relazione uscente debole, contiene un riferimento ad un oggetto DocumentDataService per utilizzare il relativo servizio;
\item\grassetto{MaaP::Client::ModelClient::Services::DocumentEditService}\\
Relazione uscente debole, contiene un riferimento ad un oggetto DocumentEditService per utilizzare il relativo servizio;
\item\grassetto{MaaP::Client::View::document}\\
Relazione entrante, interazione con il template.
\end{itemize}

\subsubsubsubsubsection{DocumentEditCtrl}
\grassetto{Nome}\\
MaaP::Client::ControllerModelView::ControllerClient::DocumentEditCtrl\\
\grassetto{Descrizione}\\
Classe che rappresenta il controller per indirizzare le richieste di visualizzazione e modifica di un document.\\
\grassetto{Utilizzo}\\
Viene utilizzata per la indirizzare le richieste di visualizzazione di una pagina di document-edit.\\
\grassetto{Relazioni con altre classi}
\begin{itemize}
\item\grassetto{MaaP::Client::ModelClient::Services::DocumentEditService}\\
Relazione uscente debole, contiene un riferimento ad un oggetto DocumentEditService per utilizzare il relativo servizio;
\item\grassetto{MaaP::Client::View::documentEdit}\\
Relazione entrante, interazione con il template.
\end{itemize}

\subsubsubsubsubsection{IndexCtrl}
\grassetto{Nome}\\
MaaP::Client::ControllerModelView::ControllerClient::IndexCtrl\\
\grassetto{Descrizione}\\
Classe che rappresenta il controller per indirizzare le richieste visualizzazione degli indici.\\
\grassetto{Utilizzo}\\
Viene utilizzata per la indirizzare le richieste di visualizzazione degli indici presenti nel database di analisi.\\
\grassetto{Relazioni con altre classi}
\begin{itemize}
\item\grassetto{MaaP::Client::ModelClient::Services::IndexService}\\
Relazione uscente debole, contiene un riferimento ad un oggetto IndexService per utilizzare il relativo servizio;
\item\grassetto{MaaP::Client::View::indexCollection}\\
Relazione entrante, interazione con il template.
\end{itemize}

\subsubsubsubsubsection{LoginCtrl}
\grassetto{Nome}\\
MaaP::Client::ControllerModelView::ControllerClient::LoginCtrl\\
\grassetto{Descrizione}\\
Classe che rappresenta il controller per indirizzare le richieste visualizzazione del form di login.\\
\grassetto{Utilizzo}\\
Viene utilizzata per la indirizzare le richieste di visualizzazione della pagina di login.\\
\grassetto{Relazioni con altre classi}
\begin{itemize}
\item\grassetto{MaaP::Client::ModelClient::Services::Authservice}\\
Relazione uscente debole, contiene un riferimento ad un oggetto Authservice per utilizzare il relativo servizio;
\item\grassetto{MaaP::Client::ModelClient::Directives::password\_check}\\
Relazione uscente debole, contiene un riferimento ad un oggetto password\_check per utilizzare il relativo servizio;
\item\grassetto{MaaP::Client::ModelClient::Directives::unique}\\
Relazione uscente debole, contiene un riferimento ad un oggetto unique per utilizzare il relativo servizio;
\item\grassetto{MaaP::Client::View::login}\\
Relazione entrante, interazione con il template.
\end{itemize}

\subsubsubsubsubsection{NavBarCtrl}
\grassetto{Nome}\\
MaaP::Client::ControllerModelView::ControllerClient::NavBarCtrl\\
\grassetto{Descrizione}\\
Classe che rappresenta il controller per indirizzare le richieste visualizzazione della barra di navigazione.\\
\grassetto{Utilizzo}\\
Viene utilizzata per la indirizzare le richieste di visualizzazione della barra di navigazione nelle pagine del progetto.\\
\grassetto{Relazioni con altre classi}
\begin{itemize}
\item\grassetto{MaaP::Client::ModelClient::Services::CollectionListService}\\
Relazione uscente debole, contiene un riferimento ad un oggetto CollectionListService per utilizzare il relativo servizio;
\item\grassetto{MaaP::Client::ModelClient::Services::LogoutService}\\
Relazione uscente debole, contiene un riferimento ad un oggetto LogoutService per utilizzare il relativo servizio;
\item\grassetto{MaaP::Client::View::Navbar}\\
Relazione entrante, interazione con il template.
\end{itemize}

\subsubsubsubsubsection{QueryCtrl}
\grassetto{Nome}\\
MaaP::Client::ControllerModelView::ControllerClient::QueryCtrl\\
\grassetto{Descrizione}\\
Classe che rappresenta il controller per indirizzare le richieste visualizzazione delle query.\\
\grassetto{Utilizzo}\\
Viene utilizzata per la indirizzare le richieste di visualizzazione della pagina contenente la lista delle query più utilizzate.\\
\grassetto{Relazioni con altre classi}
\begin{itemize}
\item\grassetto{MaaP::Client::ModelClient::Services::QuerytService}\\
Relazione uscente debole, contiene un riferimento ad un oggetto QueryService per utilizzare il relativo servizio;
\item\grassetto{MaaP::Client::ModelClient::Services::IndexService}\\
Relazione uscente debole, contiene un riferimento ad un oggetto IndexService per utilizzare il relativo servizio;
\item\grassetto{MaaP::Client::View::queryCollection}\\
Relazione entrante, interazione con il template.
\end{itemize}

\subsubsubsubsubsection{RegisterCtrl}
\grassetto{Nome}\\
MaaP::Client::ControllerModelView::ControllerClient::RegisterCtrl\\
\grassetto{Descrizione}\\
Classe che rappresenta il controller per indirizzare le richieste di registrazione.\\
\grassetto{Utilizzo}\\
Viene utilizzata per la indirizzare le richieste di registrazione di un nuovo utente.\\
\grassetto{Relazioni con altre classi}
\begin{itemize}
\item\grassetto{MaaP::Client::ModelClient::Services::RegisterService}\\
Relazione uscente debole, contiene un riferimento ad un oggetto RegisterService per utilizzare il relativo servizio;
\item\grassetto{MaaP::Client::View::register}\\
Relazione entrante, interazione con il template.
\end{itemize}

\subsubsubsubsubsection{UsersCollectionCtrl}
\grassetto{Nome}\\
MaaP::Client::ControllerModelView::ControllerClient::UsersCollectionCtrl\\
\grassetto{Descrizione}\\
Classe che rappresenta il controller per indirizzare le richieste visualizzazione degli utenti.\\
\grassetto{Utilizzo}\\
Viene utilizzata per la indirizzare le richieste di visualizzazione della pagina contenente la lista degli utenti registrati nel sistema.\\
\grassetto{Relazioni con altre classi}
\begin{itemize}
\item\grassetto{MaaP::Client::ModelClient::Services::UserCollectionService}\\
Relazione uscente debole, contiene un riferimento ad un oggetto UserCollectionService per utilizzare il relativo servizio;
\item\grassetto{MaaP::Client::ModelClient::Services::UserEditService}\\
Relazione uscente debole, contiene un riferimento ad un oggetto UserEditService per utilizzare il relativo servizio;
\item\grassetto{MaaP::Client::View::userCollection}\\
Relazione entrante, interazione con il template.
\end{itemize}

\subsubsubsubsubsection{ProfileEditCtrl}
\grassetto{Nome}\\
MaaP::Client::ControllerModelView::ControllerClient::ProfileEditCtrl\\
\grassetto{Descrizione}\\
Classe che rappresenta il controller per indirizzare le richieste di modifica dei dati di un profilo utente.\\
\grassetto{Utilizzo}\\
Viene utilizzata per la indirizzare le richieste di modifica dei dati di un profilo utente.\\
\grassetto{Relazioni con altre classi}
\begin{itemize}
\item\grassetto{MaaP::Client::ModelClient::Services::ProfileEditService}\\
Relazione uscente debole, contiene un riferimento ad un oggetto QueryService per utilizzare il relativo servizio;
\item\grassetto{MaaP::Client::View::userEdit}\\
Relazione entrante, interazione con il template.
\end{itemize}

\subsubsubsubsubsection{ProfileCtrl}
\grassetto{Nome}\\
MaaP::Client::ControllerModelView::ControllerClient::ProfileCtrl\\
\grassetto{Descrizione}\\
Classe che rappresenta il controller per indirizzare le richieste visualizzazione di un profilo utente.\\
\grassetto{Utilizzo}\\
Viene utilizzata per la indirizzare le richieste di visualizzazione dei dati di un profilo utente.\\
\grassetto{Relazioni con altre classi}
\begin{itemize}
\item\grassetto{MaaP::Client::ModelClient::Services::ProfileDataService}\\
Relazione uscente debole, contiene un riferimento ad un oggetto ProfileDataService per utilizzare il relativo servizio;
\item\grassetto{MaaP::Client::ModelClient::Services::ProfileEditService}\\
Relazione uscente debole, contiene un riferimento ad un oggetto ProfileEditService per utilizzare il relativo servizio;
\item\grassetto{MaaP::Client::View::userProfile}\\
Relazione entrante, interazione con il template.
\end{itemize}

\subsubsubsubsubsection{UserEditCtrl}
\grassetto{Nome}\\
MaaP::Client::ControllerModelView::ControllerClient::UserEditCtrl\\
\grassetto{Descrizione}\\
Classe che rappresenta il controller per indirizzare le richieste di modifica dei dati di un utente.\\
\grassetto{Utilizzo}\\
Viene utilizzata per la indirizzare le richieste di modifica dei dati di un utente da parte di un amministratore.\\
\grassetto{Relazioni con altre classi}
\begin{itemize}
\item\grassetto{MaaP::Client::ModelClient::Services::UserEditService}\\
Relazione uscente debole, contiene un riferimento ad un oggetto UserEditService per utilizzare il relativo servizio;
\item\grassetto{MaaP::Client::View::userEdit}\\
Relazione entrante, interazione con il template.
\end{itemize}

\subsubsubsubsubsection{UserCtrl}
\grassetto{Nome}\\
MaaP::Client::ControllerModelView::ControllerClient::UserCtrl\\
\grassetto{Descrizione}\\
Classe che rappresenta il controller per indirizzare le richieste di visualizzazione dei dati di un utente.\\
\grassetto{Utilizzo}\\
Viene utilizzata per la indirizzare le richieste di visualizzazione dei dati di un utente da parte di un amministratore.\\
\grassetto{Relazioni con altre classi}
\begin{itemize}
\item\grassetto{MaaP::Client::ModelClient::Services::UserDataService}\\
Relazione uscente debole, contiene un riferimento ad un oggetto UserDataService per utilizzare il relativo servizio;
\item\grassetto{MaaP::Client::ModelClient::Services::UserEditService}\\
Relazione uscente debole, contiene un riferimento ad un oggetto UserEditService per utilizzare il relativo servizio;
\item\grassetto{MaaP::Client::View::userProfile}\\
Relazione entrante, interazione con il template.
\end{itemize}

\subsubsection{MaaP::Client::ModelClient}
\subsubsubsection{Informazioni sul package}
\immagine{./Diagrammi/ComponentiModelClient}{Componente MaaP::Client::ModelClient}
\subsubsubsection{Descrizione}
Componente Model del design pattern MVVM.
\subsubsubsection{Sotto-componenti}
\begin{itemize}
\item MaaP::Client::ModelClient::Services;
\item MaaP::Client::ModelClient::Directives;
\item MaaP::Client::ModelClient::Model.
\end{itemize}

\subsubsubsection{MaaP::Client::ModelClient::Services}
\subsubsubsubsection{Informazioni sul package}
\immagine{./Diagrammi/ComponentiServices}{Componente MaaP::Client::ModelClient::Services}
\subsubsubsubsection{Descrizione}
Componente parte del ModelClient contenente i servizi per la comunicazione con il server.
\subsubsubsubsection{Classi}

\subsubsubsubsubsection{HTTP}
\grassetto{Nome}\\
MaaP::Client::ModelClient::Services::HTTP\\
\grassetto{Descrizione}\\
Classe che rappresenta il servizio di comunicazione HTTP con il server.\\
\grassetto{Utilizzo}\\
Viene utilizzata per inviare richieste HTTP al server.\\
\grassetto{Relazioni con altre classi}
\begin{itemize}
\item\grassetto{MaaP::Client::ModelClient::Services::Authservice}\\
Relazione entrante debole, interazione con il servizio di autenticazione;
\item\grassetto{MaaP::Client::ModelClient::Services::CollectionDataService}\\
Relazione entrante debole, interazione con il servizio CollectionData;
\item\grassetto{MaaP::Client::ModelClient::Services::CollectionListService}\\
Relazione entrante debole, interazione con il servizio CollectionList;
\item\grassetto{MaaP::Client::ModelClient::Services::DocumentDataService}\\
Relazione entrante debole, interazione con il servizio DocumentData;
\item\grassetto{MaaP::Client::ModelClient::Services::DocumentEditService}\\
Relazione entrante debole, interazione con il servizio DocumentEdit;
\item\grassetto{MaaP::Client::ModelClient::Services::IndexService}\\
Relazione entrante debole, interazione con il servizio degli indici;
\item\grassetto{MaaP::Client::ModelClient::Services::LogoutService}\\
Relazione entrante debole, interazione con il servizio di logout;
\item\grassetto{MaaP::Client::ModelClient::Services::ProfileDataService}\\
Relazione entrante debole, interazione con il servizio ProfileData;
\item\grassetto{MaaP::Client::ModelClient::Services::ProfileEditService}\\
Relazione entrante debole, interazione con il servizio ProfileEdit;
\item\grassetto{MaaP::Client::ModelClient::Services::QueryService}\\
Relazione entrante debole, interazione con il servizio delle query;
\item\grassetto{MaaP::Client::ModelClient::Services::RegisterService}\\
Relazione entrante debole, interazione con il servizio di registrazione;
\item\grassetto{MaaP::Client::ModelClient::Services::UserCollectionService}\\
Relazione entrante debole, interazione con il servizio UserCollection;
\item\grassetto{MaaP::Client::ModelClient::Services::UserDataService}\\
Relazione entrante debole, interazione con il servizio UserData;
\item\grassetto{MaaP::Client::ModelClient::Services::UserEditService}\\
Relazione entrante debole, interazione con il servizio UserEdit.
\end{itemize}

\subsubsubsubsubsection{Authservice}
\grassetto{Nome}\\
MaaP::Client::ModelClient::Services::Authservice\\
\grassetto{Descrizione}\\
Classe che rappresenta il servizio di gestione dell'autenticazione utente.\\
\grassetto{Utilizzo}\\
Viene utilizzata per inviare richieste di autenticazione e deautenticazione al server.\\
\grassetto{Relazioni con altre classi}
\begin{itemize}
\item\grassetto{MaaP::Client::ControllerModelView::ControllerClient::LoginCtrl}\\
Relazione entrante debole, interazione con il controller del login;
\end{itemize}

\subsubsubsubsubsection{CollectionDataService}
\grassetto{Nome}\\
MaaP::Client::ModelClient::Services::CollectionDataService\\
\grassetto{Descrizione}\\
Classe che rappresenta il servizio di gestione delle richieste dei dati di una collection.\\
\grassetto{Utilizzo}\\
Viene utilizzata per richiedere i dati di una determinata collection al server.\\
\grassetto{Relazioni con altre classi}
\begin{itemize}
\item\grassetto{MaaP::Client::ControllerModelView::ControllerClient::CollectionCtrl}\\
Relazione entrante debole, interazione con il controller delle collections;
\end{itemize}

\subsubsubsubsubsection{CollectionListService}
\grassetto{Nome}\\
MaaP::Client::ModelClient::Services::CollectionListService\\
\grassetto{Descrizione}\\
Classe che rappresenta il servizio di gestione delle richieste dei dati delle collections presenti nel sistema.\\
\grassetto{Utilizzo}\\
Viene utilizzata per richiedere al server la lista delle collections presenti nel sistema.\\
\grassetto{Relazioni con altre classi}
\begin{itemize}
\item\grassetto{MaaP::Client::ControllerModelView::ControllerClient::DashboardCtrl}\\
Relazione entrante debole, interazione con il controller della dashboard;
\item\grassetto{MaaP::Client::ControllerModelView::ControllerClient::NavBarCtrl}\\
Relazione entrante debole, interazione con il controller della barra di navigazione;
\end{itemize}

\subsubsubsubsubsection{DocumentDataService}
\grassetto{Nome}\\
MaaP::Client::ModelClient::Services::DocumentDataService\\
\grassetto{Descrizione}\\
Classe che rappresenta il servizio di gestione delle richieste dei dati dei documents presenti nel sistema.\\
\grassetto{Utilizzo}\\
Viene utilizzata per richiedere al server i dati relativi ad uno specifico document.\\
\grassetto{Relazioni con altre classi}
\begin{itemize}
\item\grassetto{MaaP::Client::ControllerModelView::ControllerClient::DocumentCtrl}\\
Relazione entrante debole, interazione con il controller dei document;
\end{itemize}

\subsubsubsubsubsection{DocumentEditService}
\grassetto{Nome}\\
MaaP::Client::ModelClient::Services::DocumentEditService\\
\grassetto{Descrizione}\\
Classe che rappresenta il servizio di gestione delle richieste dei dati dei documents presenti nel sistema.\\
\grassetto{Utilizzo}\\
Viene utilizzata per richiedere al server i dati relativi ad uno specifico document da modificare.\\
\grassetto{Relazioni con altre classi}
\begin{itemize}
\item\grassetto{MaaP::Client::ControllerModelView::ControllerClient::DocumentCtrl}\\
Relazione entrante debole, interazione con il controller dei document;
\item\grassetto{MaaP::Client::ControllerModelView::ControllerClient::DocumentEditCtrl}\\
Relazione entrante debole, interazione con il controller di modifica dei document;
\item\grassetto{MaaP::Client::ControllerModelView::ControllerClient::CollectionCtrl}\\
Relazione entrante debole, interazione con il controller delle collections;
\end{itemize}

\subsubsubsubsubsection{IndexService}
\grassetto{Nome}\\
MaaP::Client::ModelClient::Services::IndexService\\
\grassetto{Descrizione}\\
Classe che rappresenta il servizio di gestione delle richieste dei dati degli indici presenti nel database di analisi.\\
\grassetto{Utilizzo}\\
Viene utilizzata per richiedere al server i dati relativi agli indici presenti nel database di analisi.\\
\grassetto{Relazioni con altre classi}
\begin{itemize}
\item\grassetto{MaaP::Client::ControllerModelView::ControllerClient::IndexCtrl}\\
Relazione entrante debole, interazione con il controller degli indici;
\item\grassetto{MaaP::Client::ControllerModelView::ControllerClient::QueryCtrl}\\
Relazione entrante debole, interazione con il controller delle query;
\end{itemize}

\subsubsubsubsubsection{LogoutService}
\grassetto{Nome}\\
MaaP::Client::ModelClient::Services::LogoutService\\
\grassetto{Descrizione}\\
Classe che rappresenta il servizio di gestione delle richieste di logout.\\
\grassetto{Utilizzo}\\
Viene utilizzata per richiedere al server di effettuare il logout.\\
\grassetto{Relazioni con altre classi}
\begin{itemize}
\item\grassetto{MaaP::Client::ControllerModelView::ControllerClient:NavBarCtrl}\\
Relazione entrante debole, interazione con il controller della barra di navigazione;
\end{itemize}

\subsubsubsubsubsection{ProfileDataService}
\grassetto{Nome}\\
MaaP::Client::ModelClient::Services::ProfileDataService\\
\grassetto{Descrizione}\\
Classe che rappresenta il servizio di gestione delle richieste dei dati del profilo utente.\\
\grassetto{Utilizzo}\\
Viene utilizzata per richiedere al server i dati relativi ad uno specifico profilo utente.\\
\grassetto{Relazioni con altre classi}
\begin{itemize}
\item\grassetto{MaaP::Client::ControllerModelView::ControllerClient::ProfileCtrl}\\
Relazione entrante debole, interazione con il controller del profilo;
\end{itemize}

\subsubsubsubsubsection{ProfileEditService}
\grassetto{Nome}\\
MaaP::Client::ModelClient::Services::ProfileEditService\\
\grassetto{Descrizione}\\
Classe che rappresenta il servizio di gestione delle richieste dei dati del profilo utente da modificare.\\
\grassetto{Utilizzo}\\
Viene utilizzata per richiedere al server i dati relativi ad uno specifico profilo utente da modificare.\\
\grassetto{Relazioni con altre classi}
\begin{itemize}
\item\grassetto{MaaP::Client::ControllerModelView::ControllerClient::ProfileCtrl}\\
Relazione entrante debole, interazione con il controller del profilo;
\item\grassetto{MaaP::Client::ControllerModelView::ControllerClient::ProfileEditCtrl}\\
Relazione entrante debole, interazione con il controller della modifica del profilo;
\end{itemize}

\subsubsubsubsubsection{QueryService}
\grassetto{Nome}\\
MaaP::Client::ModelClient::Services::QueryService\\
\grassetto{Descrizione}\\
Classe che rappresenta il servizio di gestione delle richieste dei dati delle query.\\
\grassetto{Utilizzo}\\
Viene utilizzata per richiedere al server i dati relativi alle query più utilizzate.\\
\grassetto{Relazioni con altre classi}
\begin{itemize}
\item\grassetto{MaaP::Client::ControllerModelView::ControllerClient::QueryCtrl}\\
Relazione entrante debole, interazione con il controller delle query;
\end{itemize}

\subsubsubsubsubsection{RegisterService}
\grassetto{Nome}\\
MaaP::Client::ModelClient::Services::RegisterService\\
\grassetto{Descrizione}\\
Classe che rappresenta il servizio di gestione delle richieste di registrazione utente.\\
\grassetto{Utilizzo}\\
Viene utilizzata per richiedere al server la registrazione di un nuovo utente.\\
\grassetto{Relazioni con altre classi}
\begin{itemize}
\item\grassetto{MaaP::Client::ControllerModelView::ControllerClient::RegisterCtrl}\\
Relazione entrante debole, interazione con il controller della registrazione;
\end{itemize}

\subsubsubsubsubsection{UserCollectionService}
\grassetto{Nome}\\
MaaP::Client::ModelClient::Services::UserCollectionService\\
\grassetto{Descrizione}\\
Classe che rappresenta il servizio di gestione delle richieste dei dati utenti presenti nel sistema.\\
\grassetto{Utilizzo}\\
Viene utilizzata per richiedere al server i dati relativi a tutti i profili utenti registrati nel sistema.\\
\grassetto{Relazioni con altre classi}
\begin{itemize}
\item\grassetto{MaaP::Client::ControllerModelView::ControllerClient::UsersCollectionCtrl}\\
Relazione entrante debole, interazione con il controller della collection utenti;
\end{itemize}

\subsubsubsubsubsection{UserDataService}
\grassetto{Nome}\\
MaaP::Client::ModelClient::Services::UserDataService\\
\grassetto{Descrizione}\\
Classe che rappresenta il servizio di gestione delle richieste dei dati del profilo utente.\\
\grassetto{Utilizzo}\\
Viene utilizzata per richiedere al server i dati relativi ad uno specifico profilo utente da parte di un utente amministratore.\\
\grassetto{Relazioni con altre classi}
\begin{itemize}
\item\grassetto{MaaP::Client::ControllerModelView::ControllerClient::UsersCtrl}\\
Relazione entrante debole, interazione con il controller degli utenti;
\end{itemize}

\subsubsubsubsubsection{UserEditService}
\grassetto{Nome}\\
MaaP::Client::ModelClient::Services::UserEditService\\
\grassetto{Descrizione}\\
Classe che rappresenta il servizio di gestione delle richieste di modifica dei dati del profilo utente.\\
\grassetto{Utilizzo}\\
Viene utilizzata per richiedere al server la modifica dei dati relativi ad uno specifico profilo utente da parte di un utente amministratore.\\
\grassetto{Relazioni con altre classi}
\begin{itemize}
\item\grassetto{MaaP::Client::ControllerModelView::ControllerClient::UserCtrl}\\
Relazione entrante debole, interazione con il controller del profilo;
\item\grassetto{MaaP::Client::ControllerModelView::ControllerClient::UserEditCtrl}\\
Relazione entrante debole, interazione con il controller di modifica del profilo;
\end{itemize}

\subsubsubsection{MaaP::Client::ModelClient::Directives}
\subsubsubsubsection{Informazioni sul package}
\immagine{./Diagrammi/ComponentiDirectives}{Componente MaaP::Client::ModelClient::Directives}
\subsubsubsubsection{Descrizione}
Componente parte del ModelClient contenente le direttive per la comunicazione con il server.
\subsubsubsubsection{Classi}

\subsubsubsubsubsection{password\_check}
\grassetto{Nome}\\
MaaP::Client::ModelClient::Directives::password\_check\\
\grassetto{Descrizione}\\
Classe che rappresenta la direttiva per il controllo password doppia.\\
\grassetto{Utilizzo}\\
Viene utilizzata per controllare che le due passwords inserite dall'utente coincidano.\\
\grassetto{Relazioni con altre classi}
\begin{itemize}
\item\grassetto{MaaP::Client::ControllerModelView::ControllerClient::LoginCtrl}\\
Relazione entrante debole, interazione con il controller del login;
\end{itemize}

\subsubsubsubsubsection{unique}
\grassetto{Nome}\\
MaaP::Client::ModelClient::Directives::unique\\
\grassetto{Descrizione}\\
Classe che rappresenta la direttiva per controllare l'unicità della mail\\
\grassetto{Utilizzo}\\
Viene utilizzata per controllare l'unicità della mail inserita dall'utente in fase di registrazione.\\
\grassetto{Relazioni con altre classi}
\begin{itemize}
\item\grassetto{MaaP::Client::ControllerModelView::ControllerClient::LoginCtrl}\\
Relazione entrante debole, interazione con il controller del login;
\end{itemize}

\subsubsubsection{MaaP::Client::ModelClient::Model}
\subsubsubsubsection{Informazioni sul package}
\immagine{./Diagrammi/ComponentiModel}{Componente MaaP::Client::ModelClient::Model}
\subsubsubsubsection{Descrizione}
Componente parte del ModelClient contenente i dati di sessione.
\subsubsubsubsection{Classi}

\subsubsubsubsubsection{SessionData}
\grassetto{Nome}\\
MaaP::Client::ModelClient::Model::SessionData\\
\grassetto{Descrizione}\\
Classe che rappresenta i dati di sessione utente.\\
\grassetto{Utilizzo}\\
Viene utilizzata memorizzare i dati di sessione del client.\\
\grassetto{Relazioni con altre classi}
\begin{itemize}
\item\grassetto{MaaP::Client::ControllerModelView::ControllerClient::ControllerAutenticazione}\\
Relazione entrante debole, interazione con il controller dell'Autenticazione.
\end{itemize}

\newpage
\subsubsection{MaaPCLI}
I seguenti diagrammi rappresentano l'architettura ad alto livello del sistema di gestione del framework denominato MaaPCLI, indicando i package, le classi e le relazioni che intercorrono tra queste.

\subsubsubsection{Informazioni sul package}
\immagine{./Diagrammi/Installer}{\gloss{Diagramma delle classi} relativo alla gestione del framework da parte dell'utente \gloss{sviluppatore}}
\subsubsubsection{Descrizione}
\gloss{Namespace} globale per lo strumento di installazione e gestione del framework. Contiene le classi che gestiscono l'installazione del framework nel sistema, l'inizializzazione e gestione dello stesso.

\subsubsubsection{Classi}

\subsubsubsubsubsection{CLI}
\grassetto{Nome}\\
MaaPCLI::CLI\\
\grassetto{Descrizione}\\
\gloss{Classe} che rappresenta l'interfaccia a riga di comando.\\
\grassetto{Utilizzo}\\
Viene utilizzata dall'\gloss{utente sviluppatore} per interagire con il framework.\\
\grassetto{Relazioni con altre classi}
\begin{itemize}
\item\grassetto{MaaPCLI::Installer}\\
Relazione uscente, contiene un riferimento ad un oggetto di tipo Installer per avviare l'installazione del framework;
\item\grassetto{MaaPCLI::ProjectFacade}\\
Relazione uscente, contiene un riferimento ad un oggetto di tipo ProjectFacade per creare un nuovo progetto, clonare uno esistente oppure eliminarlo;
\item\grassetto{MaaPCLI::InstanceManager}\\
Relazione uscente, contiene un riferimento ad un oggetto di tipo InstanceManager per istanziare una \gloss{istanza} di progetto MaaP precedentemente creato.
\end{itemize}

\subsubsubsubsubsection{Installer}
\grassetto{Nome}\\
MaaPCLI::Installer\\
\grassetto{Descrizione}\\
Classe che rappresenta lo \gloss{script} di installazione del framework.\\
\grassetto{Utilizzo}\\
Viene utilizzata dall'utente sviluppatore per installare il framework e relative dipendenze nel sistema in uso.\\
\grassetto{Relazioni con altre classi}
\begin{itemize}
\item\grassetto{MaaPCLI::CLI}\\
Relazione entrante, interazione con l'interfaccia a riga di comando.
\end{itemize}

\subsubsubsubsubsection{InstanceManager}
\grassetto{Nome}\\
MaaPCLI::InstanceManager\\
\grassetto{Descrizione}\\
Classe che rappresenta lo script per l'avvio di un'istanza di un progetto esistente.\\
\grassetto{Utilizzo}\\
Viene utilizzata dall'utente sviluppatore per avviare il server caricando un determinato progetto.\\
\grassetto{Relazioni con altre classi}
\begin{itemize}
\item\grassetto{MaaPCLI::CLI}\\
Relazione entrante, interazione con l'interfaccia a riga di comando.
\end{itemize}

\subsubsubsubsubsection{ProjectFacade}
\grassetto{Nome}\\
MaaPCLI::ProjectFacade\\
\grassetto{Descrizione}\\
Classe che rappresenta la classe \gloss{Facade} nel design pattern Facade\\
\grassetto{Utilizzo}\\
Viene utilizzata dall'utente sviluppatore per interagire con il framework per la creazione di un nuovo progetto e/o per la clonazione, eliminazione di un progetto esistente.\\
\grassetto{Relazioni con altre classi}
\begin{itemize}
\item\grassetto{MaaPCLI::ProjectCreate}\\
Relazione uscente, contiene un riferimento ad un oggetto di tipo ProjectCreate per avviare la creazione di un nuovo progetto;
\item\grassetto{MaaPCLI::ProjectClone}\\
Relazione uscente, contiene un riferimento ad un oggetto di tipo ProjectClone per avviare la clonazione di un progetto esistente;
\item\grassetto{MaaPCLI::ProjectRemove}\\
Relazione uscente, contiene un riferimento ad un oggetto di tipo ProjectRemove per eliminare un progetto esistente;
\end{itemize}

\subsubsubsubsubsection{ProjectCreate}
\grassetto{Nome}\\
MaaPCLI::ProjectCreate\\
\grassetto{Descrizione}\\
Classe che rappresenta una classe del design pattern Facade.\\
\grassetto{Utilizzo}\\
Viene utilizzata dall'utente sviluppatore per avviare la creazione di un nuovo progetto.\\
\grassetto{Relazioni con altre classi}
\begin{itemize}
\item\grassetto{MaaPCLI::ProjectFacade}\\
Relazione entrante, interazione con la facciata ProjectFacade.
\end{itemize}

\subsubsubsubsubsection{ProjectClone}
\grassetto{Nome}\\
MaaPCLI::ProjectClone\\
\grassetto{Descrizione}\\
Classe che rappresenta una classe del design pattern Facade.\\
\grassetto{Utilizzo}\\
Viene utilizzata dall'utente sviluppatore per avviare la clonazione di un progetto esistente.\\
\grassetto{Relazioni con altre classi}
\begin{itemize}
\item\grassetto{MaaPCLI::ProjectFacade}\\
Relazione entrante, interazione con la facciata ProjectFacade.
\end{itemize}

\subsubsubsubsubsection{ProjectRemove}
\grassetto{Nome}\\
MaaPCLI::ProjectRemove\\
\grassetto{Descrizione}\\
Classe che rappresenta una classe del design pattern Facade.\\
\grassetto{Utilizzo}\\
Viene utilizzata dall'utente sviluppatore per eliminare un progetto esistente.\\
\grassetto{Relazioni con altre classi}
\begin{itemize}
\item\grassetto{MaaPCLI::ProjectFacade}\\
Relazione entrante, interazione con la facciata ProjectFacade.
\end{itemize}


\newpage
\section{Diagrammi di attività}
\label{attivita}
%TEMPLATE PER DIAGRAMMA ATTIVITA' o sequenza

\subsection{Utente Business}
\immagine{./Diagrammi/Activity/attUtenteBusiness}{Diagramma attività: Utente Business}
Il diagramma precedente illustra le funzionalità disponibili all'utente business.\\
Quest'ultimo, dopo aver effettuato il login, può scegliere se navigare il menu delle collection o accedere direttamente ad una di esse, o ad un document, mediante un link diretto.
Nel caso decidesse di ricercare manualmente un documento sono disponibili filtri e ordinamenti.
Può infine modificare i propri dati personali dal suo profilo.

\subsection{Utente Business Amministratore}
\immagine{./Diagrammi/Activity/attAmministratore}{Diagramma attività: Utente Business Amministratore}
Il diagramma precedente illustra le funzionalità disponibili all'utente business amministratore.\\
Quest'ultimo ha a disposizione anche tutte le funzionalità di un normale utente business, tuttavia queste ultime sono state omesse dal diagramma per evitare ridondanza e semplificare la lettura.
L'utente amministratore può gestire i profili di tutti gli utenti, gestire gli indici disponibili e modificare o cancellare i document presenti nel database.

\subsection{Utente Business Amministratore - Gestione indici}
\immagine{./Diagrammi/Activity/GestioneIndici}{Diagramma attività: Gestione Indici}
Il diagramma precedente illustra in dettaglio la gestione degli indici da parte degli Utenti Amministratori.\\
La gestione degli indici offre due possibilità: creazione e cancellazione.
Per creare un indice, l'amministratore seleziona una query tra l'enenco delle più utilizzate e ne fa un indice.
Per l'eliminazione, l'amministratore seleziona un indice esistente e lo cancella dal sistema.

\subsection{Utente Business Amministratore - Gestione document esterna}
\immagine{./Diagrammi/Activity/attAmministratoreGDIndex}{Diagramma attività: Gestione Document esterna}
Il diagramma precedente illustra in dettaglio la gestione dei document da parte degli Utenti Amministratori.\\
L'amministratore può modificare o cancellare un document mediante i pulsanti di scelta rapida posizionati accanto al nome del document oppure aprire il document per visualizzarlo e modificarlo/cancellarlo dall'interno.


\subsection{Utente Business Amministratore - Gestione document interna}
\immagine{./Diagrammi/Activity/attAmministratoreGDIndexsub}{Diagramma attività: Gestione Document interna}
Il diagramma precedente illustra in dettaglio la gestione dei document da parte degli Utenti Amministratori.\\
In questo diagramma viene descritto il caso in cui il documento viene modificato, dopo essere stato aperto, dalla sua schermata di visualizzazione.

\subsection{Utente Business Amministratore - Gestione utenti}
\immagine{./Diagrammi/Activity/attAmministratoreGU}{Diagramma attività: Gestione utenti}
Il diagramma precedente illustra la gestione degli utenti da parte degli Utenti Amministratori.\\
L'amministratore può ordinale l'elenco utenti per una chiave o utilizzando un filto. Dopo aver selezionato un profilo utente, è libero di modificarne i dati personali, i permessi o di cancellarlo dal sistema.

\subsection{Utente Sviluppatore - Gestione progetto}
\immagine{./Diagrammi/Activity/attSviluppatore}{Diagramma attività: Gestione progetto}
Il diagramma precedente illustra la gestione del progetto da parte di Utenti Sviluppatori.\\
L'Utente Sviluppatore è libero di creare, clonare o cancellare un progetto. Una volta selezionato un progetto da gestire, è libero di avviare o fermare il server MaaP, modificare i template per le pagine web e modificare od inserire nuovi file DSL, oltre che a modificare altre impostazioni minori del server.


\section{Diagrammi di sequenza}
\label{sequenza}
\subsection{Modifica della View}
\immagine{./Diagrammi/sequence}{Diagramma sequenza: Modifica View}
Il diagramma precedente illustra la sequenza di operazioni che avviene alla modifica della view da parte dell'utente.

\newpage
\section{Design Pattern}
I Design Pattern sono soluzioni a problemi ricorrenti. Adottare i Design Pattern semplifica l'attività di progettazione, rende l'architettura più \gloss{manutenibile} e favorisce il riutilizzo del codice.\\
I design pattern possono essere suddivisi in:

\begin{itemize}
\item \grassetto{Design pattern architetturali:} definiscono l'architettura dell'applicazione ad un livello più elevato;
\item \grassetto{Design pattern creazionali:} consentono di nascondere i costruttori delle classi, permettendo di creare oggetti senza conoscere la loro implementazione;
\item \grassetto{Design pattern strutturali:} consentono di riutilizzare classi pre-esistenti, fornendo un'interfaccia più adatta;
\item \grassetto{Design pattern comportamentali:} definiscono soluzioni per le interazioni tra oggetti.
\end{itemize}
Per una descrizione generale ed approfondita dei Design Pattern utilizzati si veda l'Appendice A.\\
Nella realizzazione del progetto MaaP si è deciso di implementare i seguenti Design Pattern:\\

\subsection{Design Pattern architetturali}
\subsubsection{MVVM}
%\immagine{./Diagrammi/ComponentiClient}{Applicazione di MVVM in MaaP}
\begin{itemize}
\item \grassetto{Scopo:} Il pattern MVVM è stato scelto per separare la logica dell'applicazione lato client dalla rappresentazione grafica;
\item \grassetto{Utilizzo:} Nel progetto MaaP la scelta di utilizzare AngularJS come base di partenza per l'applicazione lato client ha implicitamente comportato l'utilizzo del design pattern MVVM perché proprio di AngularJS.
\end{itemize}

\subsection{Design Pattern comportamentali}
\subsubsection{Strategy}
\immagine{./Diagrammi/DPStrategy}{Applicazione di Strategy in MaaP}
\begin{itemize}
\item \grassetto{Scopo:} Il pattern Strategy viene usato per isolare più algoritmi che svolgono la stessa funzione dal codice che esegue la funzione;
\item \grassetto{Utilizzo:} In MaaP è stato usato gestire inizialmente un singolo algoritmo di parsing del file di descrizione, ma permetterà in futuro di aggiungere nuovi algoritmi di parsing differenziati senza modificare le classi che ne fanno uso.\\
La concrete strategy attualmente presente è: DSLParser.
\end{itemize}

\subsection{Design Pattern strutturali}
\subsubsection{Adapter}
\immagine{./Diagrammi/DPAdapter}{Applicazione di Adapter in MaaP}
\begin{itemize}
\item \grassetto{Scopo:} Il pattern Adapter viene utilizzato per adattare una classe riutilizzando un oggetto già esistente. Questo semplifica l'eventuale processo di sostituzione dell'oggetto esistente, creando un'interfaccia stabile per il resto dell'applicazione;
\item \grassetto{Utilizzo:} In MaaP è stato usato per adattare la classe Passport nel Controller. PassportAdapter adatta Passport.
\end{itemize}

\subsubsection{Facade}
\immagine{./Diagrammi/DPFacade}{Applicazione di Facade in MaaP}
\begin{itemize}
\item \grassetto{Scopo:} Il pattern Facade viene usato per fornire un'interfaccia unica a più classi;
\item \grassetto{Utilizzo:} In MaaP, ProjectFacade è una Facade che presenta un'interfaccia per tutti gli oggetti gestiscono la creazione e/o modifica di un progetto:\\
\begin{itemize}
\item ProjectCreate;
\item ProjectClone;
\item ProjectRemove.
\end{itemize}
\end{itemize}

\newpage
\section{Stime di fattibilità e di bisogno di risorse}
L'architettura definita precedentemente ha raggiunto un livello di dettaglio sufficiente per fornire una stima sulla fattibilità e di bisogno delle risorse.\\
L'analisi dell'architettura progettata ha permesso di constatare che le tecnologie che si è scelto di adottare risultano sufficientemente adeguate per la realizzazione del prodotto e riescono a ricoprire le esigenze progettuali.\\
Gli strumenti scelti sono conosciuti dalla maggioranza dei componenti del gruppo che si impegneranno comunque ad approfondire le loro conoscenze inerenti all'utilizzo degli stessi.\\
Gli strumenti utilizzati sono:
\begin{itemize}
\item \grassetto{Node.js:} per la realizzazione dell'\gloss{applicazione web} lato server;
\item \grassetto{AngularJS:} per la realizzazione dell'applicazione web lato client;
\item \grassetto{Mongoose:} per la comunicazione con il database MongoDB;
\item \grassetto{Express:} framework per Node.js;
\item \grassetto{Passport:} modulo per la gestione dell'autenticazione utente;
\item \grassetto{PegJS:} generatore di parser javascript per il file di descrizione.
\end{itemize}

\section{Tracciamento componenti requisiti}
In questa sezione verranno tracciati i componenti con i requisiti. Per semplicità di lettura i requisiti associati ai figli di un componente non sono riportati anche nel padre.
\subsubsection{Tracciamento componenti - requisiti} 

\begin{center}
\begin{longtable}{|p{0.8\linewidth}|c|}
\toprule
\multicolumn{1}{|p{0.8\linewidth}}{\textbf{Componente}} & \multicolumn{1}{|c|}{\textbf{Requisito}}\\
\midrule
\endfirsthead
\multicolumn{2}{l}{\footnotesize\itshape\tablename~\thetable: continua dalla pagina precedente} \\
\toprule
\multicolumn{1}{|p{0.8\linewidth}}{\textbf{Componente}} & \multicolumn{1}{|c|}{\textbf{Requisito}}\\
\midrule
\endhead
\midrule
\multicolumn{2}{r}{\footnotesize\itshape\tablename~\thetable: continua nella prossima pagina} \\
\endfoot
\bottomrule
\caption{Tracciamento componenti - requisiti}
\label{tab:Tracciamento componenti - requisiti}\\
\endlastfoot

\midrule 
MaaPCLI
& \\

\midrule 
MaaPCLI::Installer
& \\

\midrule 
MaaPCLI::CLI
& \\

\midrule 
MaaPCLI::InstanceManager
& \\

\midrule 
MaaPCLI::ProjectFacade
& \\

\midrule 
MaaPCLI::ProjectCreate
& ROF1\\
& ROF1.1\\
& ROF1.2\\
& ROF1.3\\
& ROF1.4\\
& ROF1.4.1\\
& ROF5\\
& RDF5.3\\

\midrule 
MaaPCLI::ProjectClone
& ROF1.6\\

\midrule 
MaaPCLI::ProjectRemove
& ROF1.5\\

\midrule 
MaaP 
& \\

\midrule 
MaaP::Server
& \\

\midrule 
MaaP::Server::ModelServer::DataManager 
& \\

\midrule
MaaP::Server::ModelServer::DataManager::DatabaseAnalysisManager 
& \\

\midrule 
MaaP::Server::ModelServer::DataManager::DatabaseAnalysisManager ::DatabaseAnalysisManager
& ROF10\\
& ROF10.1\\
& ROF10.1.1\\ 
& RDF10.2\\
& RDF10.2.1\\
& RDF10.2.1.1\\
& RDF10.2.1.2\\
& RDF10.2.2\\
& RDF10.2.3\\
& ROF10.4\\
& ROF10.5\\
& ROF10.5.2\\
& ROF10.6\\
& ROF10.7\\
& ROF10.7.1.2\\
& ROF10.7.2.2\\
& ROF10.7.3\\

\midrule
MaaP::Server::ModelServer::DataManager::DatabaseAnalysisManager:: DataRetrieverAnalysis
& ROF10.6\\ 
& ROF10\\
& ROF10.1\\
& ROF10.1.1\\
& RDF10.2\\
& RDF10.2.1\\
& RDF10.2.1.1\\
& RDF10.2.1.2\\
& RDF10.2.2\\
& RDF10.2.3\\
& ROF10.4\\
& ROF10.5\\
& ROF10.5.2\\

\midrule 
MaaP::Server::ModelServer::DataManager::DatabaseUserManager 
& \\

\midrule 
MaaP::Server::ModelServer::DataManager::DatabaseUserManager:: DatabaseUserManager
& ROF10.3\\
& ROF10.3.1\\
& ROF10.3.1.2\\
& ROF10.3.1.4\\
& ROF10.3.1.5\\
& ROF10.3.2\\
& ROF10.3.3\\


\midrule 
MaaP::Server::ModelServer::DataManager::DatabaseUserManager:: DataRetrieverUsers
& ROF10.3\\
& ROF10.3.1\\
& ROF10.3.1.2\\
& ROF10.3.1.4\\
& ROF10.3.1.5\\
& ROF10.3.2\\
& ROF10.3.3\\


\midrule 
MaaP::Server::ModelServer::DataManager::IndexManager
& \\

\midrule 
MaaP::Server::ModelServer::DataManager::IndexManager::IndexManager
& ROF10.7\\
& ROF10.7.1.2\\
& ROF10.7.2.2\\
& ROF10.7.3\\

\midrule 
MaaP::Server::ModelServer::DataManager::JSonComposer
& \\

\midrule 
MaaP::Server::ModelServer::DataManager::IDatabaseManager
& ROF10\\
& ROF10.1\\
& ROF10.1.1\\
& RDF10.2\\
& RDF10.2.1\\
& RDF10.2.1.1\\
& RDF10.2.1.2\\
& RDF10.2.2\\
& RDF10.2.3\\
& ROF10.3\\
& ROF10.3.1\\
& ROF10.3.1.2\\
& ROF10.3.1.4\\
& ROF10.3.1.5\\
& ROF10.3.2\\
& ROF10.3.3\\
& ROF10.4\\
& ROF10.5\\
& ROF10.5.2\\
& ROF10.6\\
& ROF10.7\\
& ROF10.7.1.2\\
& ROF10.7.2.2\\
& ROF10.7.3\\


\midrule 
MaaP::Server::ModelServer::DataManager::IDataRetriever
& ROF10\\
& ROF10.1\\
& ROF10.1.1\\
& RDF10.2\\
& RDF10.2.1\\
& RDF10.2.1.1\\
& RDF10.2.1.2\\
& RDF10.2.2\\
& RDF10.2.3\\
& ROF10.3\\
& ROF10.3.1\\
& ROF10.3.1.2\\
& ROF10.3.1.4\\
& ROF10.3.1.5\\
& ROF10.3.2\\
& ROF10.3.3\\
& ROF10.4\\
& ROF10.5\\
& ROF10.5.2\\
& ROF10.6\\


\midrule 
MaaP::Server::ModelServer::Database
& \\

\midrule 
MaaP::Server::ModelServer::Database::MongooseDBAnalysis
& ROF10.4\\
& ROF10.5\\
& ROF10.5.2\\
& ROF10.6\\
& ROF10.7\\
& ROF10.7.1.2\\
& ROF10.7.2.2\\
& ROF10.7.3\\

\midrule 
MaaP::Server::ModelServer::Database::DBAnalysis
& ROF10.4\\
& ROF10.5\\
& ROF10.5.2\\
& ROF10.6\\
& ROF10.7\\
& ROF10.7.1.2\\
& ROF10.7.2.2\\
& ROF10.7.3\\

\midrule 
MaaP::Server::ModelServer::Database::Mongoose
& \\

\midrule 
MaaP::Server::ModelServer::Database::MongooseDBFramework
& ROF7\\
& ROF7.1\\
& ROF7.2\\
& ROF7.2.1\\
& ROF10.3\\
& ROF10.3.1\\
& ROF10.3.1.2\\
& ROF10.3.1.4\\
& ROF10.3.1.5\\
& ROF10.3.2\\
& ROF10.3.3\\


\midrule 
MaaP::Server::ModelServer::Database::DBFramework
& ROF7\\
& ROF7.1\\
& ROF7.2\\
& ROF7.2.1\\
& ROF10.3\\
& ROF10.3.1\\
& ROF10.3.1.2\\
& ROF10.3.1.4\\
& ROF10.3.1.5\\
& ROF10.3.2\\
& ROF10.3.3\\


\midrule 
MaaP::Server::ModelServer::Database::User
& ROF7\\
& ROF7.1\\
& ROF7.2\\
& ROF7.2.1\\
& ROF10.3\\
& ROF10.3.1\\
& ROF10.3.1.2\\

\midrule 
MaaP::Server::ModelServer::Database::Query
& ROF10.6\\

\midrule 
MaaP::Server::ModelServer::DSL::DSLParser
& ROF4\\

\midrule 
MaaP::Server::ModelServer::DSL::ParserInterface
& ROF4\\

\midrule 
MaaP::Server::ModelServer::DSL::DSLManager
& ROF4\\

\midrule 
MaaP::Server::ModelServer::DSL::DSLDescriptionFile
& ROF3\\

\midrule 
MaaP::Server::ModelServer::DSL::CollectionData
& \\

\midrule 
MaaP::Server::Controller 
& \\

\midrule 
MaaP::Server::Controller::Dispatcher
& \\

\midrule 
MaaP::Server::Controller::FrontController
& ROF7\\
& ROF7.1\\
& ROF7.2\\
& ROF7.2.1\\
& RDF8\\
& RDF8.1\\
& RDF8.2\\
& RDF8.2.1\\
& ROF9\\
& ROF10.2.4\\
& ROF10.3\\
& ROF10.4\\
& ROF10.5\\
& ROF10.6\\
& ROF10.7\\
& ROF10.7.1.2\\
& ROF10.7.2.2\\
& ROF10.7.3\\

\midrule 
MaaP::Server::Controller::IPassport
& ROF7\\
& ROF7.1\\
& ROF7.2\\
& ROF7.2.1\\
& ROF10.2.4\\

\midrule 
MaaP::Server::Controller::PassportAdapter 
& ROF7\\
& ROF7.1\\
& ROF7.2\\
& ROF7.2.1\\
& ROF10.2.4\\

\midrule 
MaaP::Server::Controller::Passport
& ROF7\\
& ROF7.1\\
& ROF7.2\\
& ROF7.2.1\\
& ROF10.2.4\\

\midrule 
MaaP::Client 
& \\

\midrule 
MaaP::Client::View 
& \\

\midrule 
MaaP::Client::View::Template
& \\

\midrule 
MaaP::Client::View::Template::SignIn
& ROF7\\
& ROF7.1\\
& ROF7.2\\
& ROF7.2.1\\

\midrule 
MaaP::Client::View::Template::SingUp
& RDF8\\
& RDF8.1\\
& RDF8.2\\
& RDF8.2.1\\

\midrule 
MaaP::Client::View::Template::AdminMainPageCollection
& ROF10\\
& RDF10.2\\
& RDF10.2.1\\
& RDF10.2.1.1\\
& RDF10.2.1.2\\
& RDF10.2.2\\
& RDF10.2.3\\
& ROF10.2.4\\
& ROF10.2.5\\
& ROF10.4\\
& ROF10.5\\
& ROF10.6\\


\midrule 
MaaP::Client::View::Template::UserMainPageCollection
& ROF10\\
& RDF10.2\\
& RDF10.2.1\\
& RDF10.2.1.1\\
& RDF10.2.1.2\\
& RDF10.2.2\\
& RDF10.2.3\\
& ROF10.2.4\\
& ROF10.2.5\\

\midrule 
MaaP::Client::View::Template::AdminMainPageDocument
& ROF10.1\\
& ROF10.1.1\\
& ROF10.1.2\\
& ROF10.2.4\\
& ROF10.2.5\\
& ROF10.4\\
& ROF10.5\\


\midrule 
MaaP::Client::View::Template::UserMainPageDocument
& ROF10.1\\
& ROF10.1.1\\
& ROF10.2.4\\
& ROF10.2.5\\


\midrule 
MaaP::Client::View::Template::MainPageDocumentEdit
& ROF10.1.3\\
& ROF10.2.4\\
& ROF10.2.5\\
& ROF10.5.1\\
& ROF10.5.2\\
& ROF10.5.3\\


\midrule 
MaaP::Client::View::Template::UserProfileEdit
& ROF10.2.4\\
& ROF10.2.5\\
& ROF10.3.1.1\\
& ROF10.3.1.2\\
& ROF10.3.1.3\\

\midrule 
MaaP::Client::View::Template::UserProfile
& ROF10.2.4\\
& ROF10.2.5\\
& ROF10.3\\
& ROF10.3.1\\

\midrule 
MaaP::Client::View::Template::AdminProfile
& ROF10.2.4\\
& ROF10.2.5\\
& ROF10.3.1.4\\
& ROF10.3.1.5\\
& ROF10.3.2\\
& ROF10.3.3\\


\midrule 
MaaP::Client::View::Template::PasswordRecovery
& ROF9\\
& ROF10.2.4\\

\midrule 
MaaP::Client::View::Template::IndexPage
& ROF10.6\\
& ROF10.7\\
& ROF10.7.1\\
& ROF10.7.1.1\\
& ROF10.7.1.2\\
& ROF10.7.2\\
& ROF10.7.2.1\\
& ROF10.7.2.2\\
& ROF10.7.3\\

\midrule 
MaaP::Client::ControllerModelView
& \\

\midrule 
MaaP::Client::ControllerModelView::ControllerClient
& \\

\midrule 
MaaP::Client::ControllerModelView::ControllerClient::ControllerAutenticazione
& ROF7\\
& ROF7.1\\
& ROF7.2\\
& ROF7.2.1\\
& RDF8\\
& RDF8.1\\
& RDF8.2\\
& RDF8.2.1\\

\midrule 
MaaP::Client::ControllerModelView::ControllerClient::ControllerCollection
& ROF10\\
& RDF10.2\\
& RDF10.2.1\\
& RDF10.2.1.1\\
& RDF10.2.1.2\\
& RDF10.2.2\\
& RDF10.2.3\\
& ROF10.4\\
& ROF10.5\\


\midrule 
MaaP::Client::ControllerModelView::ControllerClient::ControllerDocument
& ROF10.1\\
& ROF10.1.1\\
& ROF10.1.2\\
& ROF10.1.3\\
& ROF10.5.1\\
& ROF10.5.2\\
& ROF10.5.3\\


\midrule 
MaaP::Client::ControllerModelView::ControllerClient::ControllerProfilo
& ROF9\\
& ROF10.3\\
& ROF10.3.1\\
& ROF10.3.1.1\\
& ROF10.3.1.2\\
& ROF10.3.1.3\\
& ROF10.3.1.4\\
& ROF10.3.1.5\\
& ROF10.3.2\\
& ROF10.3.3\\


\midrule 
MaaP::Client::ControllerModelView::ControllerClient::ControllerMenu
& ROF10.2.4\\
& ROF10.2.5\\

\midrule 
MaaP::Client::ControllerModelView::ControllerClient::ControllerIndici
& ROF10.6\\
& ROF10.7\\
& ROF10.7.1\\
& ROF10.7.1.1\\
& ROF10.7.1.2\\
& ROF10.7.2\\
& ROF10.7.2.1\\
& ROF10.7.2.2\\
& ROF10.7.3\\


\midrule 
MaaP::Client::ControllerModelView::Scope
& \\

\midrule 
MaaP::Client::ControllerModelView::Scope::Collection
& ROF10\\
& RDF10.2\\
& RDF10.2.1\\
& RDF10.2.1.1\\
& RDF10.2.1.2\\
& RDF10.2.2\\
& RDF10.2.3\\
& ROF10.4\\
& ROF10.5\\

\midrule 
MaaP::Client::ControllerModelView::Scope::Query
& ROF10.6\\
& ROF10.7.1.1\\
& ROF10.7.1.2\\
& ROF10.7.3\\

\midrule 
MaaP::Client::ControllerModelView::Scope::Document
& ROF10.1\\
& ROF10.1.1\\
& ROF10.5.1\\
& ROF10.5.2\\
& ROF10.5.3\\


\midrule 
MaaP::Client::ControllerModelView::Scope::Profilo
& ROF10.3\\
& ROF10.3.1\\
& ROF10.3.1.1\\
& ROF10.3.1.2\\
& ROF10.3.1.4\\
& ROF10.3.1.5\\
& ROF10.3.2\\
& ROF10.3.3\\


\midrule 
MaaP::Client::ControllerModelView::Scope::Menu
& ROF10.2.4\\

\midrule 
MaaP::Client::ModelClient
& \\

\midrule 
MaaP::Client::ModelClient::Services
& \\

\midrule 
MaaP::Client::ModelClient::Services::HTTP
& \\

\midrule 
MaaP::Client::ModelClient::Model
& \\

\midrule 
MaaP::Client::ModelClient::Model::SessionData
& \\

%FINE TABELLA TRACCIAMENTO COMPONENTI REQUISITI
\end{longtable}
\end{center}

\subsubsection{Tracciamento requisiti - componenti}
\begin{center}
\begin{longtable}{|c|p{0.25\linewidth}|p{0.5\linewidth}|}
\toprule
\multicolumn{1}{|c|}{\textbf{Requisito}}
& \multicolumn{1}{|p{0.25\linewidth}}{\textbf{Descrizione}} 
& \multicolumn{1}{|p{0.5\linewidth}|}{\textbf{Componente}}\\
\midrule
\endfirsthead
\multicolumn{2}{l}{\footnotesize\itshape\tablename~\thetable: continua dalla pagina precedente} \\
\toprule
\multicolumn{1}{|c|}{\textbf{Requisito}}
& \multicolumn{1}{|p{0.25\linewidth}}{\textbf{Descrizione}} 
& \multicolumn{1}{|p{0.5\linewidth}|}{\textbf{Componente}}\\
\midrule
\endhead
\midrule
\multicolumn{2}{r}{\footnotesize\itshape\tablename~\thetable: continua nella prossima pagina} \\
\endfoot
\bottomrule
\caption{Tracciamento requisiti - componenti}
\label{tab:Tracciamento requisiti - componenti}\\
\endlastfoot

\midrule
ROF1
& Il sistema MaaP deve essere in grado di generare lo scheletro del progetto
& MaaPCLI::ProjectCreate\\

\midrule
ROF1.1
& Il sistema MaaP deve installare le librerie necessarie al funzionamento del progetto
& MaaPCLI::ProjectCreate\\

\midrule
ROF1.2
& Il sistema MaaP deve generare i file di configurazione necessari al funzionamento del progetto
& MaaPCLI::ProjectCreate\\

\midrule
ROF1.3
& Il sistema MaaP deve generare le directory necessarie al funzionamento del progetto
& MaaPCLI::ProjectCreate\\

\midrule
ROF1.4
& Il sistema MaaP deve generare il sistema di autenticazione per le pagine web
& MaaPCLI::ProjectCreate\\

\midrule
ROF1.4.1
& Il sistema di autenticazione per le pagine web deve essere generato insieme ad un profilo amministratore di default
& MaaPCLI::ProjectCreate\\

\midrule
ROF1.5
& Il sistema MaaP deve essere in grado eliminare un progetto esistente
& MaaPCLI::ProjectRemove\\

\midrule
ROF1.6
& Il sistema MaaP deve essere in grado di clonare un progetto esistente
& MaaPCLI::ProjectClone\\

\midrule
RFF2
& Il sistema MaaP deve permettere all'utente sviluppatore di utilizzare un editor interno specializzato per la scrittura/modifica dei file di descrizione
& \\

\midrule
RFF2.1
& Il sistema MaaP deve permettere all'utente sviluppatore di utilizzare un editor interno specializzato per la scrittura di un nuovo file di descrizione
& \\

\midrule
RFF2.1.1
& Il sistema MaaP deve permettere all'utente sviluppatore di utilizzare un editor interno specializzato per scrivere il codice del file di descrizione che intende creare
& \\

\midrule
RFF2.1.2
& Il sistema MaaP deve permettere all'utente sviluppatore di utilizzare un editor interno specializzato per salvare il codice scritto in modo permanente
& \\

\midrule
RFF2.2
& Il sistema MaaP deve permettere all'utente sviluppatore di utilizzare un editor interno specializzato per la modifica di un file di descrizione esistente
& \\

\midrule
RFF2.2.1
& Il sistema MaaP deve permettere all'utente sviluppatore di utilizzare un editor interno specializzato per modificare il codice di un file di descrizione esistente
& \\

\midrule
RFF2.2.2
& Il sistema MaaP deve permettere all'utente sviluppatore di utilizzare un editor interno specializzato per salvare il codice modificato in modo permanente
& \\

\midrule
RFF2.2.3
& Il sistema MaaP deve permettere all'utente sviluppatore di utilizzare un editor interno specializzato per annullare le modifiche al codice del file di descrizione modificato
& \\

\midrule
ROF3
& Il sistema MaaP deve permette all'utente sviluppatore di inserire un file di descrizione
& MaaP::Server::ModelServer::DSL:: DSLDescriptionFile\\

\midrule
ROF4
& Il sistema deve permette all'utente sviluppatore di utilizzare un file di descrizione
& MaaP::Server::ModelServer::DSL:: DSLParser\\
& & MaaP::Server::ModelServer::DSL:: ParserInterface\\
& & MaaP::Server::ModelServer::DSL:: DSLManager\\

\midrule
ROF4.1
& Il sistema MaaP deve permettere all'utente sviluppatore di creare la visualizzazione della Collection
& \\

\midrule
ROF4.1.1
& Il sistema MaaP deve permettere all'utente sviluppatore di creare la visualizzazione del menù per le Collection
& \\

\midrule
ROF4.1.1.1
& Il sistema MaaP deve permettere all'utente sviluppatore di definire il nome della voce relativa alla Collection
& \\

\midrule
ROF4.1.1.2
& Il sistema MaaP deve permettere all'utente sviluppatore di definire la posizione di una voce all'interno del menù
& \\

\midrule
ROF4.1.2
& Il sistema MaaP deve permettere all'utente sviluppatore di creare la visualizzazione della pagina Collection-Index
& \\

\midrule
ROF4.1.2.1
& Il sistema MaaP deve permettere all'utente sviluppatore di aggiungere delle chiavi da visualizzare nella pagina Collection-Index
& \\

\midrule
ROF4.1.2.1.1
& Il sistema MaaP deve permettere all'utente sviluppatore di aggiungere definire un'etichetta per la chiave da visualizzare
& \\

\midrule
ROF4.1.2.1.2
& Il sistema MaaP deve permettere all'utente sviluppatore di definire un campo associato alla chiave da visualizzare
& \\

\midrule
ROF4.1.2.1.3
& Il sistema MaaP deve permettere all'utente sviluppatore di definire un campo associato alla chiave da visualizzare, proveniente da un documento esterno
& \\

\midrule
ROF4.1.2.1.4
& Il sistema MaaP deve permettere all'utente sviluppatore di definire un campo associato alla chiave da visualizzare, proveniente dal risultato di una query
& \\

\midrule
ROF4.1.2.1.5
& Il sistema MaaP deve permettere all'utente sviluppatore di definire un campo associato alla chiave da visualizzare come trasformazione
& \\

\midrule
ROF4.1.2.2
& Il sistema MaaP deve permettere all'utente sviluppatore di definire un ordinamento rispetto a una chiave
& \\

\midrule
ROF4.1.2.3
& Il sistema deve permettere all'utente sviluppatore di  definire un numero massimo di Document da visualizzare per la pagina Collection-Index
& \\

\midrule
RFF4.1.2.4
& Il sistema MaaP deve permettere all'utente sviluppatore di aggiungere dei pulsanti all'interno della pagina Collection-Index
& \\

\midrule
ROF4.1.3
& Il sistema MaaP deve permettere all'utente sviluppatore di creare la visualizzazione per la pagine Document-Show
& \\

\midrule
ROF4.1.3.1
& Il sistema MaaP deve permettere all'utente sviluppatore di aggiungere delle chiavi da visualizzare nella pagina Document-Show
& \\

\midrule
ROF4.1.3.1.1
& Il sistema MaaP deve permettere all'utente sviluppatore di aggiungere definire un'etichetta per la chiave da visualizzare
& \\

\midrule
ROF4.1.3.1.2
& Il sistema MaaP deve permettere all'utente sviluppatore di definire un campo associato alla chiave da visualizzare
& \\

\midrule
ROF4.1.3.1.3
& Il sistema MaaP deve permettere all'utente sviluppatore di definire un campo associato alla chiave da visualizzare, proveniente da un documento esterno
& \\

\midrule
ROF4.1.3.1.4
& Il sistema MaaP deve permettere all'utente sviluppatore di definire un campo associato alla chiave da visualizzare, proveniente dal risultato di una query
& \\

\midrule
ROF4.1.3.1.5
& Il sistema MaaP deve permettere all'utente sviluppatore di definire un campo associato alla chiave da visualizzare come trasformazione
& \\

\midrule
RFF4.1.3.2
& Il sistema MaaP deve permettere all'utente sviluppatore di aggiungere un pulsante all'interno della  pagina Document-Show
& \\

\midrule
ROF4.2
& Il sistema MaaP deve permettere all'utente sviluppatore di modificare la visualizzazione della Collection
& \\

\midrule
ROF4.2.1
& Il sistema MaaP deve permettere all'utente sviluppatore di impostare la visualizzazione del menù delle le Collection
& \\

\midrule
ROF4.2.1.1
& Il sistema MaaP deve permettere all'utente sviluppatore di modificare il nome della voce relativa alla Collection
& \\

\midrule
ROF4.2.1.2
& Il sistema MaaP deve permettere all'utente sviluppatore di modificare la posizione di una voce all'interno del menù
& \\

\midrule
ROF4.2.2
& Il sistema MaaP deve permettere all'utente sviluppatore di impostare la visualizzazione della pagina Collection-Index
& \\

\midrule
ROF4.2.2.1
& Il sistema MaaP deve permettere all'utente sviluppatore di aggiungere delle chiavi da visualizzare nella pagina Collection-Index
& \\

\midrule
ROF4.2.2.2
& Il sistema MaaP deve permettere all'utente sviluppatore di eliminare delle chiavi da visualizzare nella pagina Collection-Index
& \\

\midrule
ROF4.2.2.3
& Il sistema MaaP deve permettere all'utente sviluppatore di definire un ordinamento, alfabetico crescente o decrescente, rispetto a una chiave
& \\

\midrule
ROF4.2.2.4
& Il sistema MaaP deve permettere all'utente sviluppatore di eliminare un ordinamento rispetto a una chiave
& \\

\midrule
ROF4.2.2.5
& Il sistema deve permettere all'utente sviluppatore di  definire un numero massimo di Document da visualizzare per la pagina Collection-Index
& \\

\midrule
ROF4.2.2.6
& Il sistema deve permettere all'utente sviluppatore di  eliminare il numero massimo di Document da visualizzare per la pagina Collection-Index
& \\

\midrule
RFF4.2.2.7
& Il sistema MaaP deve permettere all'utente sviluppatore di aggiungere dei pulsanti all'interno della pagina Collection-Index, specificando il nome del pulsante e l'azione che deve eseguire
& \\

\midrule
RFF4.2.2.8
& Il sistema MaaP deve permettere all'utente sviluppatore di eliminare dei pulsanti all'interno della pagina Collection-Index
& \\

\midrule
ROF4.2.3
& Il sistema MaaP deve permettere all'utente sviluppatore di impostare la visualizzazione per la pagine Document-Show
& \\

\midrule
ROF4.2.3.1
& Il sistema MaaP deve permettere all'utente sviluppatore di aggiungere delle chiavi da visualizzare nella pagina Document-Show
& \\

\midrule
ROF4.2.3.2
& Il sistema MaaP deve permettere all'utente sviluppatore di eliminare delle chiavi da visualizzare nella pagina Document-Show
& \\

\midrule
RFF4.2.3.3
& Il sistema MaaP deve permettere all'utente sviluppatore di aggiungere dei pulsanti all'interno della  pagina Document-Show, specificando il nome del pulsante e l'azione che deve eseguire
& \\

\midrule
RFF4.2.3.4
& Il sistema MaaP deve permettere all'utente sviluppatore di eliminare dei pulsanti all'interno della  pagina Document-Show
& \\

\midrule
ROF4.3
& Il sistema MaaP deve permettere all'utente sviluppatore di definire una query personalizzata
& \\

\midrule
ROF4.4
& Il sistema MaaP deve permettere all'utente sviluppatore di eliminare una query personalizzata
& \\

\midrule
ROF5
& Il sistema deve permettere all'utente sviluppatore la modifica dei file di configurazione
& MaaPCLI::ProjectCreate\\

\midrule
RDF5.1
& Il sistema deve permettere all'utente sviluppatore di abilitare la funzionalità di registrazione per l'utente finale nelle pagine web. Nel caso la registrazione sia abilitata l'utente finale può registrarsi al sistema, altrimenti no.
& \\

\midrule
RFF5.2
& Il sistema MaaP deve permettere all'utente sviluppatore di abilitare la funzionalità per la creazione di nuovi Document all'interno della pagina Collection-Index
& \\

\midrule
RDF5.3
& Il sistema MaaP deve permettere all'utente sviluppatore di modificare i template per le pagine web
& MaaPCLI::ProjectCreate\\

\midrule
ROF5.4
& Il sistema deve permettere all'utente sviluppatore di specificare nome, indirizzo e password, relativi al database di analisi con il quale interagire
& \\

\midrule
ROF5.5
& Il sistema deve permettere all'utente sviluppatore di abilitare la funzionalità per la creazione di nuovi indici all'interno della pagina Collection-Index
& \\


%FINE TABELLA REQUISITI MAAP, 
%Requisiti utente business

\midrule
ROF6
& L'utente business, al primo accesso, deve poter usare il profilo amministratore di default
& \\

\midrule
ROF7
& L'utente business deve potersi autenticare inserendo dei dati personali
& MaaP::Server::Controller:: FrontController\\

\midrule
ROF7.1
& L'utente business deve inserire l'email per l'autenticazione
& MaaP::Server::Controller:: FrontController\\

\midrule
ROF7.2
& L'utente business deve inserire la password per l'autenticazione
& MaaP::Server::Controller:: FrontController\\

\midrule
ROF7.2.1
& La password per l'autenticazione deve essere alfanumerica e contenere almeno otto caratteri
& MaaP::Server::Controller:: FrontController\\

\midrule
RDF8
& L'utente business deve potersi registrare inserendo dei dati personali
& MaaP::Server::Controller:: FrontController\\
& & MaaP::Client::View::Template:: SignUp\\
& & MaaP::Client::ControllerModelView:: ControllerClient::ControllerAutenticazione\\

\midrule
RDF8.1
& L'utente business, per registrarsi, deve inserire una email non presente nel sistema
& MaaP::Server::Controller:: FrontController\\
& & MaaP::Client::View::Template:: SignUp\\
& & MaaP::Client::ControllerModelView:: ControllerClient::ControllerAutenticazione\\

\midrule
RDF8.2
& L'utente business, per registrarsi, deve inserire una password
& MaaP::Server::Controller:: FrontController\\
& & MaaP::Client::View::Template:: SignUp\\
& & MaaP::Client::ControllerModelView:: ControllerClient::ControllerAutenticazione\\

\midrule
RDF8.2.1
& La password per la registrazione deve essere alfanumerica e contenere almeno otto caratteri
& MaaP::Server::Controller:: FrontController\\
& & MaaP::Client::View::Template:: SignUp\\
& & MaaP::Client::ControllerModelView:: ControllerClient::ControllerAutenticazione\\

\midrule
ROF9
& L'utente business deve poter recuperare la password
& MaaP::Server::Controller:: FrontController\\
& & MaaP::Client::View::Template:: PasswordRecovery\\
& & MaaP::Client::ControllerModelView:: ControllerClient::ControllerProfilo\\

\midrule
ROF10
& L'utente business autenticato deve poter aprire una Collection e visualizzare la sua pagina Collection-Index
& MaaP::Server::ModelServer::DataManager:: DatabaseAnalysisManager::DatabaseAnalysisManager\\
& & MaaP::Server::ModelServer::DataManager:: DatabaseAnalysisManager::DataRetrieverAnalysis\\
& & MaaP::Server::ModelServer::DataManager:: IDatabaseManager\\
& & MaaP::Server::ModelServer::DataManager:: IDataRetriever\\
& & MaaP::Client::View::Template:: AdminMainPageCollection\\
& & MaaP::Client::View::Template:: UserMainPageCollection\\
& & MaaP::Client::ControllerModelView:: ControllerClient::ControllerCollection\\
& & MaaP::Client::ControllerModelView:: Scope::Collection\\

\midrule
ROF10.1
& L'utente business autenticato deve poter visualizzare una pagina Document-Show
& MaaP::Server::ModelServer::DataManager:: DatabaseAnalysisManager::DatabaseAnalysisManager\\
& & MaaP::Server::ModelServer::DataManager:: DatabaseAnalysisManager::DataRetrieverAnalysis\\
& & MaaP::Server::ModelServer::DataManager:: IDatabaseManager\\
& & MaaP::Server::ModelServer::DataManager:: IDataRetriever\\
& & MaaP::Client::View::Template:: AdminMainPageDocument\\
& & MaaP::Client::View::Template:: UserMainPageDocument\\
& & MaaP::Client::ControllerModelView:: ControllerClient::ControllerDocument\\
& & MaaP::Client::ControllerModelView:: Scope::Document\\

\midrule
ROF10.1.1
& L'utente business autenticato deve poter visualizzare il Document selezionato
& MaaP::Server::ModelServer::DataManager:: DatabaseAnalysisManager::DatabaseAnalysisManager\\
& & MaaP::Server::ModelServer::DataManager:: DatabaseAnalysisManager::DataRetrieverAnalysis\\
& & MaaP::Server::ModelServer::DataManager:: IDatabaseManager\\
& & MaaP::Server::ModelServer::DataManager:: IDataRetriever\\
& & MaaP::Client::View::Template:: AdminMainPageDocument\\
& & MaaP::Client::View::Template:: UserMainPageDocument\\
& & MaaP::Client::ControllerModelView:: ControllerClient::ControllerDocument\\
& & MaaP::Client::ControllerModelView:: Scope::Document\\

\midrule
ROF10.1.2
& L'utente business autenticato amministratore deve poter eliminare il Document che sta visualizzando
& MaaP::Client::View::Template:: AdminMainPageDocument\\
& & MaaP::Client::ControllerModelView:: ControllerClient::ControllerDocument\\

\midrule
ROF10.1.3
& L'utente business autenticato amministratore deve poter modificare il Document che sta visualizzando
& MaaP::Client::View::Template:: MainPageDocumentEdit\\
& & MaaP::Client::ControllerModelView:: ControllerClient::ControllerDocument\\


\midrule
RDF10.2
& L'utente business autenticato deve poter modificare la visualizzazione dei Document
& MaaP::Server::ModelServer::DataManager:: DatabaseAnalysisManager::DatabaseAnalysisManager\\
& & MaaP::Server::ModelServer::DataManager:: DatabaseAnalysisManager::DataRetrieverAnalysis\\
& & MaaP::Server::ModelServer::DataManager:: IDatabaseManager\\
& & MaaP::Server::ModelServer::DataManager:: IDataRetriever\\
& & MaaP::Client::View::Template:: AdminMainPageCollection\\
& & MaaP::Client::View::Template:: UserMainPageCollection\\
& & MaaP::Client::ControllerModelView:: ControllerClient::ControllerCollection\\
& & MaaP::Client::ControllerModelView:: Scope::Collection\\

\midrule
RDF10.2.1
& L'utente business autenticato deve poter selezionare dei criteri per la visualizzazione
& MaaP::Server::ModelServer::DataManager:: DatabaseAnalysisManager::DatabaseAnalysisManager\\
& & MaaP::Server::ModelServer::DataManager:: DatabaseAnalysisManager::DataRetrieverAnalysis\\
& & MaaP::Server::ModelServer::DataManager:: IDatabaseManager\\
& & MaaP::Server::ModelServer::DataManager:: IDataRetriever\\
& & MaaP::Client::View::Template:: AdminMainPageCollection\\
& & MaaP::Client::View::Template:: UserMainPageCollection\\
& & MaaP::Client::ControllerModelView:: ControllerClient::ControllerCollection\\
& & MaaP::Client::ControllerModelView:: Scope::Collection\\

\midrule
RDF10.2.1.1
& L'utente business autenticato deve poter effettuare un ordinamento rispetto a una chiave
& MaaP::Server::ModelServer::DataManager:: DatabaseAnalysisManager::DatabaseAnalysisManager\\
& & MaaP::Server::ModelServer::DataManager:: DatabaseAnalysisManager::DataRetrieverAnalysis\\
& & MaaP::Server::ModelServer::DataManager:: IDatabaseManager\\
& & MaaP::Server::ModelServer::DataManager:: IDataRetriever\\
& & MaaP::Client::View::Template:: AdminMainPageCollection\\
& & MaaP::Client::View::Template:: UserMainPageCollection\\
& & MaaP::Client::ControllerModelView:: ControllerClient::ControllerCollection\\
& & MaaP::Client::ControllerModelView:: Scope::Collection\\

\midrule
RDF10.2.1.2
& L'utente business deve poter selezionare un numero massimo di Document da visualizzare per pagina
& MaaP::Server::ModelServer::DataManager:: DatabaseAnalysisManager::DatabaseAnalysisManager\\
& & MaaP::Server::ModelServer::DataManager:: DatabaseAnalysisManager::DataRetrieverAnalysis\\
& & MaaP::Server::ModelServer::DataManager:: IDatabaseManager\\
& & MaaP::Server::ModelServer::DataManager:: IDataRetriever\\
& & MaaP::Client::View::Template:: AdminMainPageCollection\\
& & MaaP::Client::View::Template:: UserMainPageCollection\\
& & MaaP::Client::ControllerModelView:: ControllerClient::ControllerCollection\\
& & MaaP::Client::ControllerModelView:: Scope::Collection\\

\midrule
RDF10.2.2
& L'utente business autenticato deve poter applicare un filtro alla visualizzazione dei Document
& MaaP::Server::ModelServer::DataManager:: DatabaseAnalysisManager::DatabaseAnalysisManager\\
& & MaaP::Server::ModelServer::DataManager:: DatabaseAnalysisManager::DataRetrieverAnalysis\\
& & MaaP::Server::ModelServer::DataManager:: IDatabaseManager\\
& & MaaP::Server::ModelServer::DataManager:: IDataRetriever\\
& & MaaP::Client::View::Template:: AdminMainPageCollection\\
& & MaaP::Client::View::Template:: UserMainPageCollection\\
& & MaaP::Client::ControllerModelView:: ControllerClient::ControllerCollection\\
& & MaaP::Client::ControllerModelView:: Scope::Collection\\

\midrule
RDF10.2.3
& L'utente business autenticato deve poter annullare il filtro
& MaaP::Server::ModelServer::DataManager:: DatabaseAnalysisManager::DatabaseAnalysisManager\\
& & MaaP::Server::ModelServer::DataManager:: DatabaseAnalysisManager::DataRetrieverAnalysis\\
& & MaaP::Server::ModelServer::DataManager:: IDatabaseManager\\
& & MaaP::Server::ModelServer::DataManager:: IDataRetriever\\
& & MaaP::Client::View::Template:: AdminMainPageCollection\\
& & MaaP::Client::View::Template:: UserMainPageCollection\\
& & MaaP::Client::ControllerModelView:: ControllerClient::ControllerCollection\\
& & MaaP::Client::ControllerModelView:: Scope::Collection\\

\midrule
ROF10.2.4
& L'utente business autenticato deve poter disconnettersi
& MaaP::Server::Controller:: FrontController\\
& & MaaP::Server::Controller:: IPassport\\
& & MaaP::Server::Controller:: PassportAdapter\\
& & MaaP::Server::Controller:: Passport\\
& & MaaP::Client::View::Template:: AdminMainPageCollection\\
& & MaaP::Client::View::Template:: UserMainPageCollection\\
& & MaaP::Client::View::Template:: AdminMainPageDocument\\
& & MaaP::Client::View::Template:: UserMainPageDocument\\
& & MaaP::Client::View::Template:: MainPageDocumentEdit\\
& & MaaP::Client::View::Template:: UserProfileEdit\\
& & MaaP::Client::View::Template:: UserProfile\\
& & MaaP::Client::View::Template:: AdminProfile\\
& & MaaP::Client::View::Template:: PasswordRecovery\\
& & MaaP::Client::ControllerModelView:: ControllerClient::ControllerMenu\\
& & MaaP::Client::ControllerModelView:: Scope::Menu\\

\midrule
ROF10.2.5
& L'utente business autenticato deve poter navigare tra la Collection
& MaaP::Client::View::Template:: AdminMainPageCollection\\
& & MaaP::Client::View::Template:: UserMainPageCollection\\
& & MaaP::Client::View::Template:: AdminMainPageDocument\\
& & MaaP::Client::View::Template:: UserMainPageDocument\\
& & MaaP::Client::View::Template:: MainPageDocumentEdit\\
& & MaaP::Client::View::Template:: UserProfileEdit\\
& & MaaP::Client::View::Template:: UserProfile\\
& & MaaP::Client::View::Template:: AdminProfile\\
& & MaaP::Client::ControllerModelView:: ControllerClient::ControllerMenu\\

\midrule
ROF10.3
& L'utente business autenticato deve poter gestire il proprio profilo
& MaaP::Server::ModelServer::DataManager:: DatabaseUserManager::DatabaseUserManager\\
& & MaaP::Server::ModelServer::DataManager:: DatabaseUserManager::DataRetrieverUsers\\
& & MaaP::Server::ModelServer::DataManager:: IDatabaseManager\\
& & MaaP::Server::ModelServer::DataManager:: IDataRetriever\\
& & MaaP::Server::ModelServer::Database:: MongooseDBFramework\\
& & MaaP::Server::ModelServer::Database:: DBFramework\\
& & MaaP::Server::ModelServer::Database:: User\\
& & MaaP::Server::Controller:: FrontController\\
& & MaaP::Client::View::Template:: UserProfile\\
& & MaaP::Client::ControllerModelView:: ControllerClient::ControllerProfilo\\
& & MaaP::Client::ControllerModelView:: Scope::Profilo\\

\midrule
ROF10.3.1
& L'utente business autenticato deve poter gestire i propri dati
& MaaP::Server::ModelServer::DataManager:: DatabaseUserManager::DatabaseUserManager\\
& & MaaP::Server::ModelServer::DataManager:: DatabaseUserManager::DataRetrieverUsers\\
& & MaaP::Server::ModelServer::DataManager:: IDatabaseManager\\
& & MaaP::Server::ModelServer::DataManager:: IDataRetriever\\
& & MaaP::Server::ModelServer::Database:: MongooseDBFramework\\
& & MaaP::Server::ModelServer::Database:: DBFramework\\
& & MaaP::Server::ModelServer::Database:: User\\
& & MaaP::Client::View::Template:: UserProfile\\
& & MaaP::Client::ControllerModelView:: ControllerClient::ControllerProfilo\\
& & MaaP::Client::ControllerModelView:: Scope::Profilo\\

\midrule
ROF10.3.1.1
& L'utente business autenticato deve poter modificare i propri dati utente
& MaaP::Client::View::Template:: UserProfileEdit\\
& & MaaP::Client::ControllerModelView:: ControllerClient::ControllerProfilo\\
& & MaaP::Client::ControllerModelView:: Scope::Profilo\\

\midrule
ROF10.3.1.2
& L'utente business autenticato deve poter  salvare le modifiche apportate
& MaaP::Server::ModelServer::DataManager:: DatabaseUserManager::DatabaseUserManager\\
& & MaaP::Server::ModelServer::DataManager:: DatabaseUserManager::DataRetrieverUsers\\
& & MaaP::Server::ModelServer::DataManager:: IDatabaseManager\\
& & MaaP::Server::ModelServer::DataManager:: IDataRetriever\\
& & MaaP::Server::ModelServer::Database:: MongooseDBFramework\\
& & MaaP::Server::ModelServer::Database:: DBFramework\\
& & MaaP::Server::ModelServer::Database:: User\\
& & MaaP::Client::View::Template:: UserProfileEdit\\
& & MaaP::Client::ControllerModelView:: ControllerClient::ControllerProfilo\\
& & MaaP::Client::ControllerModelView:: Scope::Profilo\\

\midrule
ROF10.3.1.3
& L'utente business autenticato deve poter annullare le modifiche apportate
& MaaP::Client::View::Template:: UserProfileEdit\\
& & MaaP::Client::ControllerModelView:: ControllerClient::ControllerProfilo\\

\midrule
ROF10.3.1.4
& L'utente business autenticato amministratore deve poter modificare i dati degli utenti business
& MaaP::Server::ModelServer::DataManager:: DatabaseUserManager::DatabaseUserManager\\
& & MaaP::Server::ModelServer::DataManager:: DatabaseUserManager::DataRetrieverUsers\\
& & MaaP::Server::ModelServer::DataManager:: IDatabaseManager\\
& & MaaP::Server::ModelServer::DataManager:: IDataRetriever\\
& & MaaP::Server::ModelServer::Database:: MongooseDBFramework\\
& & MaaP::Server::ModelServer::Database:: DBFramework\\
& & MaaP::Client::View::Template:: AdminProfile\\
& & MaaP::Client::ControllerModelView:: ControllerClient::ControllerProfilo\\
& & MaaP::Client::ControllerModelView:: Scope::Profilo\\

\midrule
ROF10.3.1.5
& L'utente business autenticato deve poter modificare i permessi degli utenti business
& MaaP::Server::ModelServer::DataManager:: DatabaseUserManager::DatabaseUserManager\\
& & MaaP::Server::ModelServer::DataManager:: DatabaseUserManager::DataRetrieverUsers\\
& & MaaP::Server::ModelServer::DataManager:: IDatabaseManager\\
& & MaaP::Server::ModelServer::DataManager:: IDataRetriever\\
& & MaaP::Server::ModelServer::Database:: MongooseDBFramework\\
& & MaaP::Server::ModelServer::Database:: DBFramework\\
& & MaaP::Client::View::Template:: AdminProfile\\
& & MaaP::Client::ControllerModelView:: ControllerClient::ControllerProfilo\\
& & MaaP::Client::ControllerModelView:: Scope::Profilo\\

\midrule
ROF10.3.2
& L'utente business autenticato amministratore deve poter creare un nuovo utente business
& MaaP::Server::ModelServer::DataManager:: DatabaseUserManager::DatabaseUserManager\\
& & MaaP::Server::ModelServer::DataManager:: DatabaseUserManager::DataRetrieverUsers\\
& & MaaP::Server::ModelServer::DataManager:: IDatabaseManager\\
& & MaaP::Server::ModelServer::DataManager:: IDataRetriever\\
& & MaaP::Server::ModelServer::Database:: MongooseDBFramework\\
& & MaaP::Server::ModelServer::Database:: DBFramework\\
& & MaaP::Client::View::Template:: AdminProfile\\
& & MaaP::Client::ControllerModelView:: ControllerClient::ControllerProfilo\\
& & MaaP::Client::ControllerModelView:: Scope::Profilo\\

\midrule
ROF10.3.3
& L'utente business autenticato amministratore deve poter eliminare un utente business
& MaaP::Server::ModelServer::DataManager:: DatabaseUserManager::DatabaseUserManager\\
& & MaaP::Server::ModelServer::DataManager:: DatabaseUserManager::DataRetrieverUsers\\
& & MaaP::Server::ModelServer::DataManager:: IDatabaseManager\\
& & MaaP::Server::ModelServer::DataManager:: IDataRetriever\\
& & MaaP::Server::ModelServer::Database:: MongooseDBFramework\\
& & MaaP::Server::ModelServer::Database:: DBFramework\\
& & MaaP::Client::View::Template:: AdminProfile\\
& & MaaP::Client::ControllerModelView:: ControllerClient::ControllerProfilo\\
& & MaaP::Client::ControllerModelView:: Scope::Profilo\\

\midrule
ROF10.4
& L'utente business autenticato amministratore deve poter cancellare un Document
& MaaP::Server::ModelServer::DataManager:: DatabaseAnalysisManager::DatabaseAnalysisManager\\
& & MaaP::Server::ModelServer::DataManager:: DatabaseAnalysisManager::DataRetrieverAnalysis\\
& & MaaP::Server::ModelServer::DataManager:: IDatabaseManager\\
& & MaaP::Server::ModelServer::DataManager:: IDataRetriever\\
& & MaaP::Server::ModelServer::Database:: MongooseDBAnalysis\\
& & MaaP::Server::ModelServer::Database:: DBAnalysis\\
& & MaaP::Server::Controller:: FrontController\\
& & MaaP::Client::View::Template:: AdminMainPageCollection\\
& & MaaP::Client::View::Template:: AdminMainPageDocument\\
& & MaaP::Client::ControllerModelView:: ControllerClient::ControllerCollection\\
& & MaaP::Client::ControllerModelView:: Scope::Collection\\

\midrule
ROF10.5
& L'utente business autenticato amministratore deve poter modificare un Document
& MaaP::Server::ModelServer::DataManager:: DatabaseAnalysisManager::DatabaseAnalysisManager\\
& & MaaP::Server::ModelServer::DataManager:: DatabaseAnalysisManager::DataRetrieverAnalysis\\
& & MaaP::Server::ModelServer::DataManager:: IDatabaseManager\\
& & MaaP::Server::ModelServer::DataManager:: IDataRetriever\\
& & MaaP::Server::ModelServer::Database:: MongooseDBAnalysis\\
& & MaaP::Server::ModelServer::Database:: DBAnalysis\\
& & MaaP::Server::Controller:: FrontController\\
& & MaaP::Client::View::Template:: AdminMainPageCollection\\
& & MaaP::Client::View::Template:: AdminMainPageDocument\\
& & MaaP::Client::ControllerModelView:: ControllerClient::ControllerCollection\\
& & MaaP::Client::ControllerModelView:: Scope::Collection\\

\midrule
ROF10.5.1
& L'utente business autenticato amministratore deve poter modificare i valori associati alla chiavi
& MaaP::Client::View::Template:: MainPageDocumentEdit\\
& & MaaP::Client::ControllerModelView:: ControllerClient::ControllerDocument\\
& & MaaP::Client::ControllerModelView:: Scope::Document\\

\midrule
ROF10.5.2
& L'utente business autenticato amministratore deve poter salvare le modifiche apportate al Document
& MaaP::Server::ModelServer::DataManager:: DatabaseAnalysisManager::DatabaseAnalysisManager\\
& & MaaP::Server::ModelServer::DataManager:: DatabaseAnalysisManager::DataRetrieverAnalysis\\
& & MaaP::Server::ModelServer::DataManager:: IDatabaseManager\\
& & MaaP::Server::ModelServer::DataManager:: IDataRetriever\\
& & MaaP::Server::ModelServer::Database:: MongooseDBAnalysis\\
& & MaaP::Server::ModelServer::Database:: DBAnalysis\\
& & MaaP::Client::View::Template:: MainPageDocumentEdit\\
& & MaaP::Client::ControllerModelView:: ControllerClient::ControllerDocument\\
& & MaaP::Client::ControllerModelView:: Scope::Document\\

\midrule
ROF10.5.3
& L'utente business autenticato amministratore deve poter annullare le modifiche apportate al Document
& MaaP::Client::View::Template:: MainPageDocumentEdit\\
& & MaaP::Client::ControllerModelView:: ControllerClient::ControllerDocument\\
& & MaaP::Client::ControllerModelView:: Scope::Document\\

\midrule
ROF10.6
& L'utente business autenticato amministratore deve poter visualizzare le query più utilizzate dal sistema MaaP
& MaaP::Server::ModelServer::DataManager:: DatabaseAnalysisManager::DatabaseAnalysisManager\\
& & MaaP::Server::ModelServer::DataManager:: DatabaseAnalysisManager::DataRetrieverAnalysis\\
& & MaaP::Server::ModelServer::DataManager:: IDatabaseManager\\
& & MaaP::Server::ModelServer::DataManager:: IDataRetriever\\
& & MaaP::Server::ModelServer::Database:: MongooseDBAnalysis\\
& & MaaP::Server::ModelServer::Database:: DBAnalysis\\
& & MaaP::Server::ModelServer::Database:: Query\\
& & MaaP::Server::Controller:: FrontController\\
& & MaaP::Client::View::Template:: AdminMainPageCollection\\
& & MaaP::Client::View::Template:: IndexPage\\
& & MaaP::Client::ControllerModelView:: ControllerClient::ControllerIndici\\
& & MaaP::Client::ControllerModelView:: Scope::Query\\

\midrule
ROF10.7
& L'utente business autenticato amministratore deve poter gestire la creazione e l'eliminazione degli indici
& MaaP::Server::ModelServer::DataManager:: DatabaseAnalysisManager::DatabaseAnalysisManager\\
& & MaaP::Server::ModelServer::DataManager:: IndexManager::IndexManager\\
& & MaaP::Server::ModelServer::DataManager:: IDatabaseManager\\
& & MaaP::Server::ModelServer::Database:: MongooseDBAnalysis\\
& & MaaP::Server::ModelServer::Database:: DBAnalysis\\
& & MaaP::Server::ModelServer::Database:: MongooseDBFramework\\
& & MaaP::Server::ModelServer::Database:: DBFramework\\
& & MaaP::Server::ModelServer::Database:: User\\
& & MaaP::Server::Controller:: FrontController\\
& & MaaP::Server::Controller:: IPassport\\
& & MaaP::Server::Controller:: PassportAdapter\\
& & MaaP::Server::Controller:: Passport\\
& & MaaP::Client::View::Template:: SignIn\\
& & MaaP::Client::View::Template:: IndexPage\\
& & MaaP::Client::ControllerModelView:: ControllerClient::ControllerAutenticazione\\
& & MaaP::Client::ControllerModelView:: ControllerClient::ControllerIndici\\

\midrule
ROF10.7.1
& L'utente business autenticato amministratore deve poter creare degli indici
& MaaP::Server::ModelServer::Database:: MongooseDBFramework\\
& & MaaP::Server::ModelServer::Database:: DBFramework\\
& & MaaP::Server::ModelServer::Database:: User\\
& & MaaP::Server::Controller:: IPassport\\
& & MaaP::Server::Controller:: PassportAdapter\\
& & MaaP::Server::Controller:: Passport\\
& & MaaP::Client::View::Template:: SignIn\\
& & MaaP::Client::View::Template:: IndexPage\\
& & MaaP::Client::ControllerModelView:: ControllerClient::ControllerAutenticazione\\
& & MaaP::Client::ControllerModelView:: ControllerClient::ControllerIndici\\

\midrule
ROF10.7.1.1
& L'utente business autenticato amministratore deve poter selezionare una query
& MaaP::Client::View::Template:: IndexPage\\
& & MaaP::Client::ControllerModelView:: ControllerClient::ControllerIndici\\
& & MaaP::Client::ControllerModelView:: Scope::Query\\

\midrule
ROF10.7.1.2
& L'utente business autenticato amministratore deve poter creare l'indice
& MaaP::Server::ModelServer::DataManager:: DatabaseAnalysisManager::DatabaseAnalysisManager\\
& & MaaP::Server::ModelServer::DataManager:: IndexManager::IndexManager\\
& & MaaP::Server::ModelServer::DataManager:: IDatabaseManager\\
& & MaaP::Server::ModelServer::Database:: MongooseDBAnalysis\\
& & MaaP::Server::ModelServer::Database:: DBAnalysis\\
& & MaaP::Server::Controller:: FrontController\\
& & MaaP::Client::View::Template:: IndexPage\\
& & MaaP::Client::ControllerModelView:: ControllerClient::ControllerIndici\\
& & MaaP::Client::ControllerModelView:: Scope::Query\\

\midrule
ROF10.7.2
& L'utente business autenticato amministratore deve poter eliminare degli indici
& MaaP::Server::ModelServer::Database:: MongooseDBFramework\\
& & MaaP::Server::ModelServer::Database:: DBFramework\\
& & MaaP::Server::ModelServer::Database:: User\\
& & MaaP::Server::Controller:: IPassport\\
& & MaaP::Server::Controller:: PassportAdapter\\
& & MaaP::Server::Controller:: Passport\\
& & MaaP::Client::View::Template:: SignIn\\
& & MaaP::Client::View::Template:: IndexPage\\
& & MaaP::Client::ControllerModelView:: ControllerClient::ControllerAutenticazione\\
& & MaaP::Client::ControllerModelView:: ControllerClient::ControllerIndici\\

\midrule
ROF10.7.2.1
& L'utente business autenticato amministratore deve poter selezionare un indice presente nel sistema
& MaaP::Server::ModelServer::Database:: MongooseDBFramework\\
& & MaaP::Server::ModelServer::Database:: DBFramework\\
& & MaaP::Server::ModelServer::Database:: User\\
& & MaaP::Server::Controller:: IPassport\\
& & MaaP::Server::Controller:: PassportAdapter\\
& & MaaP::Server::Controller:: Passport\\
& & MaaP::Client::View::Template:: SignIn\\
& & MaaP::Client::View::Template:: IndexPage\\
& & MaaP::Client::ControllerModelView:: ControllerClient::ControllerAutenticazione\\
& & MaaP::Client::ControllerModelView:: ControllerClient::ControllerIndici\\

\midrule
ROF10.7.2.2
& L'utente business autenticato amministratore deve poter eliminare l'indice selezionato
& MaaP::Server::ModelServer::DataManager:: DatabaseAnalysisManager::DatabaseAnalysisManager\\
& & MaaP::Server::ModelServer::DataManager:: IndexManager::IndexManager\\
& & MaaP::Server::ModelServer::DataManager:: IDatabaseManager\\
& & MaaP::Server::ModelServer::Database:: MongooseDBAnalysis\\
& & MaaP::Server::ModelServer::Database:: DBAnalysis\\
& & MaaP::Server::Controller:: FrontController\\
& & MaaP::Client::View::Template:: IndexPage\\
& & MaaP::Client::ControllerModelView:: ControllerClient::ControllerIndici\\

\midrule
ROF10.7.3
& L'utente business autenticato amministratore deve poter visualizzare gli indici presenti nel sistema
& MaaP::Server::ModelServer::DataManager:: DatabaseAnalysisManager::DatabaseAnalysisManager\\
& & MaaP::Server::ModelServer::DataManager:: IndexManager::IndexManager\\
& & MaaP::Server::ModelServer::DataManager:: IDatabaseManager\\
& & MaaP::Server::ModelServer::Database:: MongooseDBAnalysis\\
& & MaaP::Server::ModelServer::Database:: DBAnalysis\\
& & MaaP::Server::Controller:: FrontController\\
& & MaaP::Client::View::Template:: IndexPage\\
& & MaaP::Client::ControllerModelView:: ControllerClient::ControllerIndici\\
& & MaaP::Client::ControllerModelView:: Scope::Query\\

%FINE TABELLA TRACCIAMENTO REQUISITI COMPONENTI
\bottomrule
\caption{Tracciamento requisiti - componenti}
\end{longtable}
\end{center} 


\newpage
\appendix
\section{Descrizione Design Pattern} %A.0

\subsection{Design Pattern architetturali}
\subsubsection{MVVM}
\immagine{./Diagrammi/MVVM}{Diagramma del design pattern MVVM}
\begin{itemize}
\item \grassetto{Scopo:} Disaccoppiare le tre componenti seguenti:
\begin{itemize}
\item Model: dati di business e regole di accesso;
\item View: rappresentazione grafica;
\item ViewModel: punto d'incontro tra View e Model. I dati ricevuti da quest'ultimo sono elaborati per essere presentati e passati alla View.
\end{itemize}
\item \grassetto{Motivazione:} Lo scopo di molte applicazioni è quello di recuperare dati e visualizzarli in maniera opportuna a seconda delle esigenze degli utenti. Poiché il flusso chiave di informazione avviene tra il dispositivo su cui sono memorizzati i dati e l'interfaccia utente, si è portati a legare insieme queste due parti per ridurre la quantità di codice e migliorare le performance dell'applicazione. Questo approccio, apparentemente naturale, presente alcuni problemi significativi; uno di questi è che l'interfaccia utente tende a cambiare più in fretta rispetto al sistema di memorizzazione dei dati. Un altro problema, che si ha nel mettere insieme i dati e l'interfaccia utente, è che le applicazioni aziendali tendono ad incorporare logica di business che va al di là della semplice trasmissione dei dati. C'è la necessità, quindi, di rendere modulari le funzionalità dell'interfaccia utente in maniera tale da poter facilmente modificare le singole parti. La soluzione a tutto ciò è costituita dal pattern Model-View-ViewModel (MVVM) che separa la modellazione del dominio, la presentazione e le azioni basate sugli input degli utenti all'interno di tre classi separate;
\item \grassetto{Applicabilità:} Il pattern MVVM può essere utilizzo nei seguenti casi:
\begin{itemize}
\item Quando si vuole trattare un gruppo di oggetti come un oggetto singolo;
\item Quando si vuole disaccoppiare View e Model instaurando un \gloss{protocollo di sottoscrizione} e notifica tra loro;
\item Quando si vogliono agganciare più View ad un Model per fornire più rappresentazioni del Model stesso.
\end{itemize}
\end{itemize}

\subsection{Design Pattern strutturali}
\subsubsection{Adapter}
\immagine{./Diagrammi/ADAPTER}{Diagramma del design pattern Adapter}
\begin{itemize}
\item \grassetto{Scopo:} Convertire l'interfaccia di una classe in un'altra interfaccia richiesta dal client. Consente a classi diverse di operare insieme quando ciò non sarebbe altrimenti possibile a causa di interfacce incompatibili;
\item \grassetto{Motivazione:} A volte una classe di supporto, che è stata progettata con obbiettivi di riuso, non può essere riusata semplicemente perché la sua interfaccia non è compatibile con l'interfaccia richiesta da un'applicazione;
\item \grassetto{Applicabilità:} Il pattern Adapter può essere utilizzato nei seguenti casi:
\begin{itemize}
\item Quando si vuole usare una classe esistente, ma la sua interfaccia non è compatibile con quella desiderata;
\item Quando si vuole creare una classe riusabile in grado di cooperare con classi non correlate o impreviste, cioè con classi che non necessariamente hanno interfacce compatibili;
\item Per gli oggetti adapter quando si devono utilizzare diverse sottoclassi esistenti, ma non è pratico adattare la loro interfaccia creando una sottoclasse per ciascuna di esse.
\end{itemize}
\end{itemize}

\subsubsection{Facade}
\immagine{./Diagrammi/FACADE}{Diagramma del design pattern Facade}
\begin{itemize}
\item \grassetto{Scopo:} Fornire un'interfaccia unificata per un insieme di interfacce presenti in un \gloss{sottosistema}. Definisce un'interfaccia di livello più alto che rende il sottosistema più semplice da utilizzare;
\item \grassetto{Motivazione:} Suddividere un sistema in sottosistemi aiuta a ridurre la complessità. Un obbiettivo comune di progettazione è la minimizzazione delle comunicazioni e delle dipendenze fra i diversi sottosistemi. Un modo per raggiungere questo obbiettivo è introdurre un oggetto facade, che fornisce un'interfaccia unica e semplificata per accedere alle funzionalità offerte da un sottosistema;
\item \grassetto{Applicabilità:} Il pattern Facade può essere utilizzato nei seguenti casi:
\begin{itemize}
\item Quando si vuole fornire un'interfaccia semplice a un sottosistema complesso poiché fornisce una vista semplice di base su un sottosistema che si rivela essere sufficiente per la maggior parte dei client;
\item Nei casi in cui si cono molte dipendenze fra i client e le classi che implementano un'astrazione in quanto si disaccoppia il sottosistema dai client e dagli altri sistemi, promuovendo \gloss{portabilità} ed indipendenza dei sottosistemi;
\item Quando si vogliono organizzare i sottosistemi in una struttura a livelli.
\end{itemize}
\end{itemize}

\subsection{Design Pattern comportamentali} %A.4
\subsubsection{Strategy}
\immagine{./Diagrammi/STRATEGY}{Diagramma del design pattern Strategy}
\begin{itemize}
\item \grassetto{Scopo:} Definire una famiglia di algoritmi, incapsularli e renderli intercambiabili. Permette agli algoritmi di variare indipendentemente dai client che ne fanno uso;
\item \grassetto{Motivazione:} Esistono molti algoritmi per risolvere un problema. Codificare statisticamente ognuno di questi algoritmi nelle classi che ne fanno richiesta non è auspicabile per svariati motivi. Si possono evitare questi problemi definendo delle classi che incapsulano svariati algoritmi chiamati Strategy;
\item \grassetto{Applicabilità:} Il pattern Strategy può essere utilizzato nei seguenti casi:
\begin{itemize}
\item Molte classi correlate differiscono fra loro solo per il comportamento;
\item Sono necessarie più varianti di un algoritmo;
\item Un algoritmo usa una \gloss{struttura dati} che non dovrebbe essere resa nota ai client;
\item Una classe definisce molti comportamenti che compaiono all'interno delle scelte condizionali multiple.
\end{itemize}
\end{itemize}


%FINE DOCUMENTO NON CANCELLARE
\end{document}
